
%\documentclass[useAMS,usenatbib,prd,preprint]{revtex4-1}
%\documentclass[a4paper, amsfonts, amssymb, amsmath, reprint, showkeys, nofootinbib, twoside, onecolumn,notitlepage]{revtex4-1}
\documentclass[aps,prl,singlecolumn,nofootinbib]{revtex4-2}


\def\aj{{AJ}}                   % Astronomical Journal

\def\actaa{{Acta Astron.}}      % Acta Astronomica
\def\araa{{ARA\&A}}             % Annual Review of Astron and Astrophys
\def\apj{{ApJ}}                 % Astrophysical Journal
\def\apjl{{ApJ}}                % Astrophysical Journal, Letters
\def\apjs{{ApJS}}               % Astrophysical Journal, Supplement
\def\ao{{Appl.~Opt.}}           % Applied Optics
\def\apss{{Ap\&SS}}             % Astrophysics and Space Science
\def\aap{{A\&A}}                % Astronomy and Astrophysics
\def\aapr{{A\&A~Rev.}}          % Astronomy and Astrophysics Reviews
\def\aaps{{A\&AS}}              % Astronomy and Astrophysics, Supplement
\def\azh{{AZh}}                 % Astronomicheskii Zhurnal
\def\baas{{BAAS}}               % Bulletin of the AAS
\def\bac{{Bull. astr. Inst. Czechosl.}}
                % Bulletin of the Astronomical Institutes of Czechoslovakia 
\def\caa{{Chinese Astron. Astrophys.}}
                % Chinese Astronomy and Astrophysics
\def\cjaa{{Chinese J. Astron. Astrophys.}}
                % Chinese Journal of Astronomy and Astrophysics
\def\icarus{{Icarus}}           % Icarus
\def\jcap{{J. Cosmology Astropart. Phys.}}
                % Journal of Cosmology and Astroparticle Physics
\def\jrasc{{JRASC}}             % Journal of the RAS of Canada
\def\memras{{MmRAS}}            % Memoirs of the RAS
\def\mnras{{MNRAS}}             % Monthly Notices of the RAS
\def\na{{New A}}                % New Astronomy
\def\nar{{New A Rev.}}          % New Astronomy Review
\def\pra{{Phys.~Rev.~A}}        % Physical Review A: General Physics
\def\prb{{Phys.~Rev.~B}}        % Physical Review B: Solid State
\def\prc{{Phys.~Rev.~C}}        % Physical Review C
\def\prd{{Phys.~Rev.~D}}        % Physical Review D
\def\pre{{Phys.~Rev.~E}}        % Physical Review E
\def\prl{{Phys.~Rev.~Lett.}}    % Physical Review Letters
\def\pasa{{PASA}}               % Publications of the Astron. Soc. of Australia
\def\pasp{{PASP}}               % Publications of the ASP
\def\pasj{{PASJ}}               % Publications of the ASJ
\def\rmxaa{{Rev. Mexicana Astron. Astrofis.}}%
                % Revista Mexicana de Astronomia y Astrofisica
\def\qjras{{QJRAS}}             % Quarterly Journal of the RAS
\def\skytel{{S\&T}}             % Sky and Telescope
\def\solphys{{Sol.~Phys.}}      % Solar Physics
\def\sovast{{Soviet~Ast.}}      % Soviet Astronomy
\def\ssr{{Space~Sci.~Rev.}}     % Space Science Reviews
\def\zap{{ZAp}}                 % Zeitschrift fuer Astrophysik
\def\nat{{Nature}}              % Nature
\def\iaucirc{{IAU~Circ.}}       % IAU Cirulars
\def\aplett{{Astrophys.~Lett.}} % Astrophysics Letters
\def\apspr{{Astrophys.~Space~Phys.~Res.}}
                % Astrophysics Space Physics Research
\def\bain{{Bull.~Astron.~Inst.~Netherlands}} 
                % Bulletin Astronomical Institute of the Netherlands
\def\fcp{{Fund.~Cosmic~Phys.}}  % Fundamental Cosmic Physics
\def\gca{{Geochim.~Cosmochim.~Acta}}   % Geochimica Cosmochimica Acta
\def\grl{{Geophys.~Res.~Lett.}} % Geophysics Research Letters
\def\jcp{{J.~Chem.~Phys.}}      % Journal of Chemical Physics
\def\jgr{{J.~Geophys.~Res.}}    % Journal of Geophysics Research
\def\jqsrt{{J.~Quant.~Spec.~Radiat.~Transf.}}
                % Journal of Quantitiative Spectroscopy and Radiative Transfer
\def\memsai{{Mem.~Soc.~Astron.~Italiana}}
                % Mem. Societa Astronomica Italiana
\def\nphysa{{Nucl.~Phys.~A}}   % Nuclear Physics A
\def\physrep{{Phys.~Rep.}}   % Physics Reports
\def\physscr{{Phys.~Scr}}   % Physica Scripta
\def\planss{{Planet.~Space~Sci.}}   % Planetary Space Science
\def\procspie{{Proc.~SPIE}}   % Proceedings of the SPIE


\bibliographystyle{apsrev4-1}



\usepackage{amsmath,amstext}
\usepackage[T1]{fontenc}
%\usepackage{apjfonts} 
%\usepackage[figure,figure*]{hypcap}

%\usepackage{amsbsy}
\usepackage{amssymb}
%\usepackage{amsmath}
%\input{epsf}
\usepackage{graphicx}
\usepackage{ae,aecompl}


\usepackage{hyperref}
\usepackage{amsmath}
\usepackage{amssymb}
\usepackage{mathtools}
\usepackage{bm}
\usepackage{cleveref}
\usepackage{tensor}
\usepackage{braket}
\usepackage{enumitem}
\usepackage{mhchem}
\usepackage{cancel}
\usepackage{amsthm}
\usepackage{upgreek}
\usepackage{multirow}


\usepackage{color}
\newcommand{\todo}[1]{{\textcolor{blue}{\texttt{TODO: #1}} }}
\newcommand{\red}[1]{{\color{red} #1}}
\newcommand{\green}[1]{{\color{green} #1}}
\newcommand{\blue}[1]{{\color{blue} #1}}
\newcommand{\vect}[1]{\mathbf{{#1}}}
\newcommand{\parallelsum}{\mathbin{\!/\mkern-5mu/\!}}
\newcommand{\spc}{\quad \quad \quad}
\newcommand{\paral}{\mathbin{\!/\mkern-5mu/\!}}
\newcommand{\facciatabianca}{\newpage\shipout\null\stepcounter{page}}
\newcommand{\riga}{\textcolor{white}{-}}

\newcommand{\ra}[1]{\renewcommand{\arraystretch}{#1}} 

\newcommand{\n}{{\mathrm{n}}}
\newcommand{\p}{{\mathrm{p}}}
\def\deq{\hat{\delta}}

\def\be{\begin{equation}}
\def\ee{\end{equation}}
\def\beq{\begin{eqnarray}}
\def\eeq{\end{eqnarray}}


\theoremstyle{definition}
\newtheorem{definition}{Definition}

\theoremstyle{theorem}
\newtheorem{theorem}{Theorem}


%Marcelo's additions
\usepackage{soul}
\setstcolor{blue}
\newcommand{\md}[1]{{\color{blue}#1}}
\newcommand{\mdNote}[1]{{\color{blue}[#1]}}

\begin{document}
\title{Universality Classes of Relativistic Fluid Dynamics: Foundations}
\author{L.~Gavassino$^1$, M.~Disconzi$^{1}$, \& J.~Noronha$^2$}
\affiliation{
$^1$Department of Mathematics, Vanderbilt University, Nashville, TN, USA
\\
$^2$Illinois Center for Advanced Studies of the Universe \& Department of Physics,
University of Illinois at Urbana-Champaign, Urbana, IL 61801-3003, USA
}


\begin{abstract}
A general organizing principle is proposed that can be used to derive the equations of motion describing the near-equilibrium dynamics of causal and thermodynamically stable relativistic systems. The latter are found to display some new type of universal behavior near equilibrium that allows them to be grouped into universality classes defined by their degrees of freedom, information content, and conservation laws. The universality classes expose a number of surprising equivalences between different theories, shedding new light on the near-equilibrium behavior of relativistic systems. 
\end{abstract} 
%
%
%Main message: It is possible to group relativistic hydrodynamic theories into ``universality classes'', where two theories belong to the same class if they are equivalent in the linear regime. Our result: In the presence of an information current, two theories belong to the same universality class if an only if there is a change of variables that maps the information current and the entropy production rate of the first theory into the information current and the entropy production rate of the second theory.
%
%Why is this beautiful?
%\begin{itemize}
%\item It explains why so many theories are the same in the linear regime.
%\item Once you know the conditions for linear stability and causality of one theory in the class, you know them for all theories in the class (see examples).
%\item You can group even substances into universality classes (see examples). 
%\item You can prove fun facts. For example: superfluid hydro coincides with Israel-Stewart heat conduction, when we send the heat conductivity to infinity.
%\end{itemize}
%
\maketitle
%\noindent
%\textit{Introduction}: 
\section{Introduction}
The construction of a relativistic hydrodynamic theory usually relies on two choices. First, one must choose the degrees of freedom that describe the hydrodynamic state. Then, a guiding principle is used to derive the dynamical equations for those fields. Such guiding principle may be, e.g., a thermodynamic principle \cite{Israel_Stewart_1979,Muller_book,Otting1_1998,Salazar2020,GavassinoFronntiers2021}, a variational principle \cite{noto_rel,Dubosky2012_effective,TorrieriIS2016,GavassinoRadiazione}, kinetic theory \cite{Denicol2012Boltzmann}, holography \cite{Bhattacharyya:2007vjd,Baier2008}, or effective theory arguments \cite{Bemfica2017TheFirst,Kovtun2019,Bemfica:2019knx,Hoult:2020eho,Bemfica-Disconzi-Graber-2021,BemficaDNDefinitivo2020}. The level of freedom involved in these two choices is immense \cite{Geroch_Lindblom_1991_causal}, leading to a plethora of alternative theories  
%(described by a variety of acronyms)
, which are constantly being added to the ``relativistic hydrodynamics landscape'' \cite{Israel_Stewart_1979,Baier2008,Denicol2012Boltzmann,Florkowski:2010cf,Martinez:2010sc,Jaiswal:2013vta,Bemfica2017TheFirst,Kovtun2019,Bemfica:2019knx,Hoult:2020eho,BemficaDNDefinitivo2020,Florkowski_2018,
RauWasserman2020,Kiamari_2021,SperanzaChiral2021,
NoronhaGeneralFrame2021,Perna_2021,
GavassinoKhalatnikov2022,KeYin2022,
Singh2022,Most_2022,
HellerSingulant2022,Brito_2022,Wagner:2022ayd,SalazarZannias2022}. 
%. This puts the predictive power of relativistic hydrodynamics as a whole into question. In fact, if we can construct (at least in principle) an infinite number of different theories to model the same system, and such theories lead to different predictions, then relativistic hydrodynamics may be effectively non-falsifiable \cite{Geroch2001}. This issue is becoming particularly pressing, now that so many new theories are being added to the ``hydrodynamic landscape'' \cite{Florkowski_2018,Kovtun2019,RauWasserman2020,Kiamari_2021,SperanzaChiral2021,NoronhaGeneralFrame2021,Perna_2021,GavassinoKhalatnikov2022,KeYin2022,Singh2022,Most_2022,HellerSingulant2022,Brito_2022}. 
Nevertheless, it is known that some theories can become equivalent in certain regimes. For instance, second-order theories for viscous hydrodynamics arising from very different choices of both fields and dynamical equations  \cite{Hishcock1983,OlsonLifsh1990,GerochLindblom1990,
carter1991,Baier2008,Denicol2012Boltzmann,
Stricker2019} are mathematically equivalent near equilibrium \cite{Liu1986,PriouCOMPAR1991,GavassinoNonHydro2022,GavassinoGENERIC2022}. Equivalency here means that, although those theories can be different in the nonlinear regime, they become indistinguishable when linearized around homogeneous equilibrium states. Demonstrating such equivalency sometimes requires making a complicated change of variables so that, almost ``magically''  (see, e.g.,  \cite{GavassinoGENERIC2022}), one linearized theory is transformed into the other.
The fact that such transformations are possible near equilibrium cannot be some fortuitous coincidence. Rather, this should follow as a consequence of the underlying properties of equilibrium states in relativity.

Here, we present a general organizing principle that can be used to derive the equations of motion (EOM) describing the near-equilibrium dynamics of any causal and thermodynamically stable relativistic system. The method is based on a new result, rigorously proven here, which establishes that the EOM describing the system's linearized disturbances can always be obtained from a 4-vector field $E^\mu$ known as the ``information current'' \cite{GavassinoCausality2021}, which is a quadratic function of the perturbations (e.g., temperature variations). Assuming an isotropic and homogeneous equilibrium state, we show how the information current can be systematically constructed for arbitrary theories in terms of the corresponding linear-order perturbation fields, grouped according to their transformation properties under the $SO(3)$ rotation group. Entropy production follows from $\partial_\mu E^\mu \leq 0$,  which describes the fact that our initial information about the microstate of the system \cite{Jaynes1965} is erased as all microstates evolve towards the equilibrium state. Furthermore, we prove that, no matter how complicated a hydrodynamic theory is, one can always rearrange its equations of motion so that they resemble Israel-Stewart theory \cite{Israel_Stewart_1979} near equilibrium. This is used to reveal that causal and thermodynamically stable theories of relativistic fluid dynamics possess a new type of universal behavior near equilibrium that allows them to be grouped into universality classes. Each class describes a physically different behavior, defined solely by the degrees of freedom, the corresponding information current, and the conservation laws. The existence of such universality classes unveils a number of startling equivalences between seemingly different sets of equations of motion. 
%Finally, in this paper we use our approach to determine the equations of motion that describe the near-equilibrium dynamics of elastic systems in relativity. 
Finally, as an application, we use our approach to show that (in the linear regime) an isotropic solid may be viewed as a fluid with an additional conserved charge, which transforms as a symmetric $(0,2)$-tensor under $SO(3)$.
Many more specific examples can be worked out for different universality classes (see accompanying paper). \emph{Notation}: We use $\hbar=k_B = c=1$, a $(-,+,+,+)$ Minkowski metric, Greek indices run from 0 to 3,  lowercase Latin indices from 1 to 3. Uppercase Latin indices are multi-index labels. $M_{(AB)}$ and $M_{[AB]}$ denote symmetric and antisymmetric parts of $M_{AB}$.

%If we want to save hydrodynamics, there is only one hope: Perhaps, in a certain regime of physical interest, most of these apparently different theories become indistinguishable, and fall into a relatively limited number of ``universality classes'', where  

%Interestingly, we already have an astonishing example of such phenomenon: Second-order theories for viscous hydrodynamic that arise from very different choices of both fields and dynamical equations  \cite{Hishcock1983,OlsonLifsh1990,GerochLindblom1990,carter1991,Baier2008,Denicol2012Boltzmann,Stricker2019} turn out to be all mathematically equivalent close to equilibrium \cite{Liu1986,PriouCOMPAR1991,GavassinoNonHydro2022,GavassinoGENERIC2022}. By ``equivalent'', we mean that the equations, although profoundly different in the non-linear regime, become literally indistinguishable (also in the UV limit!), when we linearise them around homogeneous equilibrium states. To show this, sometimes one needs to make a complicated change of variables so that, almost ``magically''  (see e.g. \cite{GavassinoGENERIC2022}), one linearised theory is transformed into another one, with all pieces in the proper place. 
%The goal of this Letter is to rigorously prove that this phenomenon is not a coincidence, but a manifestation of a general mechanism, by which most non-linear theories must fall into universality classes when linearised.

% In a nutshell, our result is the following. Take all possible linear theories which describe the evolution of a relativistic fluid for small deviations around global equilibrium (which is assumed homogeneous and isotropic), and which are consistent with the second law of thermodynamics (at least in the linear regime). Classify such theories depending on how their fields $\varphi^A$ transform under rotations in the global equilibrium rest frame. That is, count the number $\mathfrak{s}$ of rotational scalars, the number $\mathfrak{v}$ of rotational vectors, the number $\mathfrak{t}$ of rotational symmetric traceless 2tensors, and so on. For example, if one of the degrees of freedom of the perturbation is a fourcurrent $\delta n^\mu$, then $\delta n^0$ should be counted among the scalars, and $\delta n^j$ ($j=1,2,3$) should be counted among the vectors. In this Letter, we will show that all linear theories whose fields have the same geometric character, $(\mathfrak{s},\mathfrak{v},\mathfrak{t},...)$, are indeed the same linear theory, possibly written using different variables. As consequence, each choice of $(\mathfrak{s},\mathfrak{v},\mathfrak{t},...)$ defines a universality class of relativistic fluid dynamics (note, however, that some transport coefficients may be constrained to zero in a specific theory, while being different from zero in the most general theory of the class).
% Below we provide a list of the most common (parity conserving \cite{SperanzaChiral2021,GavassinoNonHydro2022}) universality classes of the form $(\mathfrak{s},\mathfrak{v},\mathfrak{t},0,0,0...)$, together with the best known representatives of each class. 

%\emph{Notation}: Throughout the paper, we adopt the following conventions and notation: ...

%\noindent
%\textit{Hyperbolicity from thermodynamics}:
\section{Hyperbolicity from thermodynamics}
To understand how the universality classes come about, we need to derive some new results concerning the properties of the EOM of relativistic systems near equilibrium. We first recap some general facts about relativistic thermodynamics. All states of thermodynamic equilibrium maximize some thermodynamic potential $\Phi$ \cite{Stuekelberg1962,Israel_2009_inbook,Grmela1997,
huang_book,landau5,GavassinoTermometri}. For equations linearized about equilibrium, the quantity $E:=\Phi_\text{eq}-\Phi \geq 0$ plays the role of a non-increasing Lyapunov functional \cite{GavassinoGibbs2021}, which can be expressed as a volume integral $E(\Sigma) = \int_\Sigma E^\mu \, d\Sigma_\mu$, 
where $\Sigma$ is a Cauchy surface, $d\Sigma_\mu$ is its normal surface element (with orientation $d\Sigma_0>0$ \cite{MTW_book}), and $E^\mu$ is the information current \cite{GavassinoCausality2021} (see the Supplemental Material for a brief review). For the state of thermodynamic equilibrium to be stable against perturbations in all reference frames \cite{GavassinoSuperluminal2021,GavassinoBounds2023}, $E^\mu$ must be future-directed and timelike for any non-vanishing perturbation \cite{GavassinoCausality2021}. The second law of thermodynamics also requires 
\begin{equation}\label{Emu}
\partial_\mu E^\mu = -\sigma \leq 0 \, ,
\end{equation}
where $\sigma$ is the entropy production rate. We consider here the dynamics of linear deviations about thermodynamic equilibrium. Hence, let $\varphi^A(t,x^j)$  be some real linear-order perturbation fields that we use to characterize the system's state. Thus, $\varphi^A=0$ in equilibrium.  We now prove the following result: 
% \md{
% \begin{theorem}\label{theo}
% Consider the following system of PDEs in $\mathbb{R}^{1+3}$,
% \begin{equation}\label{prima}
% M^\mu_{AB}\partial_\mu \varphi^B = -\Xi_{AB} \varphi^B \, ,
% \end{equation}
% where $M^\mu_{AB}$ and $\Xi_{AB}$ are constant
% matrices and $M^0_{AB}$ is invertible. Assume 
% the following. (i) 
% There exist constant 
% symmetric matrices 
% $E^\mu_{AB}$ and $\sigma_{AB}$ such that
% \eqref{Emu} holds for any smooth solution of \eqref{prima}, 
% where $E^\mu$ and $\sigma$ are given by
% \begin{equation}\label{second}
% \begin{split}
%  E^\mu ={}& \dfrac{1}{2} E^\mu_{AB} \varphi^A \varphi^B \, ,  \\
%  \sigma ={}&  \sigma_{AB} \varphi^A \varphi^B \, . \\
% \end{split}
% \end{equation}
% (ii) For any vector $Z^A$, $E_{AB} Z^A Z^B$
% is future-directed timelike.
% (iii) Given arbitrary smooth initial data $\varphi^A(0,x^j)$
% for the system \eqref{prima}, there always exists a smooth 
% solution $\varphi^A$ taking the given data.
% Then, (I) equations \eqref{prima} are causal;
% (II) if $\varphi^A$ is a smooth solution to  
% \eqref{prima}, then it is also a solution to 
% \begin{equation}\label{cinque}
% E^\mu_{AB} \partial_\mu \varphi^B = -\sigma_{AB} \varphi^B - \Xi_{[AB]}\varphi^B \, . 
% \end{equation}
% \end{theorem}
% We postpone the proof until the end of the manuscript. Here, we make some important
% remarks about Theorem \ref{theo}.
% }
\begin{theorem}\label{theo}
Consider the following system of partial differential equations on $\mathbb{R}^{1+3}$,
\begin{equation}\label{prima}
M^\mu_{AB}\partial_\mu \varphi^B = -\Xi_{AB} \varphi^B \, ,
\end{equation}
where $M^\mu_{AB}$ and $\Xi_{AB}$ are constant
matrices, and $M^0_{AB}$ is invertible. Suppose that there exist constant symmetric matrices $E^\mu_{AB}$ and $\sigma_{AB}$ such that equation \eqref{Emu} holds for any smooth solution of \eqref{prima}, 
where $E^\mu$ and $\sigma$ are given by
\begin{equation}\label{second}
 E^\mu= \dfrac{1}{2} E^\mu_{AB} \varphi^A \varphi^B\,, \, \qquad  
 \sigma =  \sigma_{AB} \varphi^A \varphi^B \,,
\end{equation}
and $E^\mu$ is future-directed timelike (and hence non-vanishing) over the support of $\varphi^A$. Then, the system \eqref{prima} is causal, and it can be equivalently rewritten in a symmetric hyperbolic form as follows:
\begin{equation}\label{cinque}
E^\mu_{AB} \partial_\mu \varphi^B = -\sigma_{AB} \varphi^B - \Xi_{[AB]}\varphi^B \,. 
\end{equation} 
\end{theorem}


% We postpone the proof until the end of the manuscript. Here, we make some important
% remarks about Theorem \ref{theo}.


% Suppose that the initial data $\varphi^A(0,x^j)$ can be specified freely (no constraints \footnote{In the presence of constraints, the theorem still applies, but one needs to redefine the fields $\varphi^A$ in such a way that they are all independent. The standard example is the shear stress tensor, which must be symmetric traceless, $\Pi^{[jk]}=\Pi\indices{^j _j}=0$. We cannot set $\varphi^A=\Pi^{jk}$, because the field components would not be independent, but we can construct $5$ independent fields $\varphi^A$ ($A=1,...,5$) such that $\Pi^{jk}=\Pi^{jk}(\varphi^A)$. See \cite{GavassinoCasmir2022} for more details.}), and that the linearised field equations are of first order,
 
%  Furthermore, suppose that equation 
 
%  which are taken symmetric (with no loss of generality). Then, if the theory is thermodynamically stable, the system \eqref{prima} can be recast in a symmetric-hyperbolic form as follows:

% Furthermore, the theory is causal. 
\begin{proof}
System \eqref{prima} admits smooth solutions of the form
\begin{equation}\label{thevictory}
    \varphi^A(t,x^j)=  \big(e^{-(M^0)^{-1} \Xi t} \big)\indices{^A_B} \, (Z^B-W^B) 
    + \big(e^{-(M^0)^{-1} (\Xi+M^j a_j) t} \big)\indices{^A_B} \, W^Be^{a_j x^j} \, , 
\end{equation}
for any real $Z^A,W^A,a_j=\text{const}$. Thus, $E^\mu_{AB}Z^AZ^B= 2E^\mu(0)$ must be timelike future directed for any $Z^A \neq 0$. In particular, $E^0_{AB}$ must be positive definite and, hence, invertible. Now, we always have the freedom to redefine the matrices $M^\mu_{AB}$ and $\Xi_{AB}$ by contracting both sides of \eqref{prima} with an invertible matrix $\mathcal{N}^A_C$. Let us then ``fix'' the matrix $M^0_{AB}$ to coincide with $E^0_{AB}$ (this is possible because both are invertible). If we contract \eqref{prima} with $\varphi^A$, and we plug \eqref{second} into \eqref{Emu}, we obtain the two equations below:
\begin{equation}\label{gringo}
\varphi^A M^\mu_{AB}\partial_\mu \varphi^B + \varphi^A \Xi_{AB} \varphi^B =0 \, , \spc  \varphi^A E^\mu_{AB} \partial_\mu \varphi^B + \varphi^A \sigma_{AB} \varphi^B =0 \, . 
\end{equation}
Both are respected along \textit{all} solutions of \eqref{prima}. If we subtract the second equation of \eqref{gringo} to the first, the terms $\varphi^A M^0_{AB}\partial_t \varphi^B$ cancel out (we have fixed $M^0_{AB}=E^0_{AB}$). Evaluating the result along \eqref{thevictory}, at $x^\mu=0$, we obtain
\begin{equation}\label{MjjjEjjj}
Z^A(M^j_{AB}-E^j_{AB})a_j W^B +Z^A (\Xi_{AB}- \sigma_{AB} ) Z^B =0 \, .
\end{equation}
Since this must be true for any choice of $Z^A$, $W^A$, and $a_j$, we obtain $M^j_{AB}=E^j_{AB}$ and $\Xi_{(AB)}= \sigma_{AB}$. We have recovered \eqref{cinque}. But the matrices $E^\mu_{AB}$ are symmetric \footnote{Note that, if in \eqref{second} we did not define $E^\mu_{AB}$ to be symmetric, in the second equation of \eqref{gringo} there would be $E^\mu_{(AB)}$ in place of $E^\mu_{AB}$, and we would eventually get $M^\mu_{AB}=E^\mu_{(AB)}$. Hence, the system would anyway be symmetric.}, and $E^0_{AB}$ is positive definite. Thus, the system \eqref{cinque} is symmetric hyperbolic and, since $E^\mu_{AB}Z^AZ^B$ is future-directed timelike for any $Z^A \neq 0$, it is also causal \cite{Geroch_Lindblom_1991_causal}.
\end{proof}
This theorem  
shows that thermodynamic stability in relativistic systems implies not only causality \cite{GavassinoCausality2021} but also symmetric hyperbolicity in the linear regime (provided that we have an information current). This is convenient given that, in most physical systems, symmetric-hyperbolicity follows directly from Onsager symmetry \cite{GavassinoCasmir2022}. Furthermore, by showing that the EOM are symmetric hyperbolic, we establish  that the initial value problem of all thermodynamically consistent theories 
linearized about homogeneous equilibrium
is well-posed   \cite{GerochPartial1996,Kato2009,CourantHilbert2_book}, i.e., given initial data, solutions to the equations always exist, are unique, and depend continuously on the data. 
The assumptions in Theorem \ref{theo} are quite general, encompassing an astounding number of different theories. For example, they are satisfied by all Israel-Stewart-like theories \cite{Israel_Stewart_1979,Hishcock1983,Baier2008,DMNR2012}, in an arbitrary hydrodynamic frame \cite{OlsonLifsh1990,BritoStability2020,
NoronhaGeneralFrame2021}, and with an arbitrary number of chemical species \cite{Almaalol2022}. They are also satisfied within Carter's multifluid theory \cite{carter1991,noto_rel,
GavassinoStabilityCarter2022}, GENERIC-based theories \cite{Otting2_1999,Otting3_1999,Stricker2019,
GavassinoGENERIC2022}, Geroch-Lindblom theories \cite{Geroch_Lindblom_1991_causal,LindblomRelaxation1996,
Geroch1995,GavassinoNonHydro2022}, and divergence-type theories \cite{Liu1986,GerochLindblom1990,
ZanniasDivergenceType2023}. Theorem \ref{theo} shows that the EOM of all of those theories can be found (and written in symmetric-hyperbolic form) in terms of the information current. Later in this paper we show how the information current can be systematically determined from symmetry arguments. In the Supplemental Material, we explore the origin of Theorem \ref{theo} in detail.    


Finally, we remark that the hypotheses of Theorem \ref{theo} are violated by first-order theories \cite{Bemfica2017TheFirst,
Bemfica2019_conformal1,Bemfica-Disconzi-Graber-2021,
BemficaDNDefinitivo2020,Kovtun2019,GavassinoLyapunov_2020,DoreGavassino2022} because their regularized information current contains derivatives \cite{GavassinoInfoBDNK2024ufs}. Also, we note that in most hydrodynamic theories, apart from the ``holographic theories'' discussed in \cite{Heller2014,GavassinoNonHydro2022}, one finds $\Xi_{[AB]}=0$. Consequently, in this case, the linear field equations \eqref{cinque} are uniquely determined by the information current and the entropy production rate. We will assume this in the following and work with theories specified by the triplet $\{\varphi^A,E^\mu,\sigma \}$. 

In this context,  one can find conditions under which two seemingly different theories are, in reality, different manifestations of the same near-equilibrium physics. This is the content of:  
\begin{theorem}\label{theo2}
Let $\{\varphi^A, E^\mu,\sigma \}$ and $\{\tilde{\varphi}^C, \tilde{E}^\mu,\tilde{\sigma} \}$ be two linear theories, for which all the hypotheses of Theorem \ref{theo} hold, and such that $\Xi_{[AB]}=\tilde{\Xi}_{[CD]}=0$. Then, such theories are equivalent if and only if there is an invertible matrix $\mathcal{N}^A_C$ such that, for arbitrary $Z^C$,
\begin{equation}\label{potato}
E^\mu(\mathcal{N}^A_C Z^C)=\tilde{E}^\mu(Z^C) \,,\qquad   \sigma(\mathcal{N}^A_C Z^C)= \tilde{\sigma}(Z^C) \,.
\end{equation}
\end{theorem}
\begin{proof}
Theorem \ref{theo} implies that the field equations of the two theories can be recast as
\begin{equation}\label{cinquantaAcinque}
E^\mu_{AB} \partial_\mu \varphi^B = -\sigma_{AB} \varphi^B \, ,\qquad 
\tilde{E}^\mu_{CD} \partial_\mu \tilde{\varphi}^D = -\tilde{\sigma}_{CD} \tilde{\varphi}^D \, .
\end{equation}
Suppose that an invertible matrix $\mathcal{N}^A_C$ that satisfies \eqref{potato} exists. Then $E^\mu_{AB}\mathcal{N}^A_C \mathcal{N}^B_D=\tilde{E}^\mu_{CD}$, and $\sigma_{AB}\mathcal{N}^A_C \mathcal{N}^B_D=\tilde{\sigma}_{CD}$. But this implies that if we contract the first equation of \eqref{cinquantaAcinque} with $\mathcal{N}^A_C$, and we make the replacement $\varphi^B=\mathcal{N}^B_D \tilde{\varphi}^D$, we obtain the second equation of \eqref{cinquantaAcinque}. Hence, the equations of the two theories are the same, just written using different variables. Vice versa, suppose that the two linear theories are the same theory. Then, since the information current is unique \cite{GavassinoCausality2021}, we must have that $E^\mu = \tilde{E}^\mu$, and $\sigma=\tilde{\sigma}$. But this is equivalent to saying that there is a one-to-one mapping $\varphi^A =\mathcal{N}^A_C \tilde{\varphi}^C$ for which equation \eqref{potato} holds, with $Z^C=\tilde{\varphi}^C$.
\end{proof}

Theorem \ref{theo2} can be employed to prove that many apparently different theories currently in use reduce to exactly the same theory close to equilibrium. For example, a fluid mixture of two chemical substances undergoing a chemical reaction is indistinguishable from the Israel-Stewart theory for bulk viscosity \cite{Causality_bulk}, close to equilibrium (see also \cite{BulkGavassino,Camelio2022,CamelioSimulations2022}). Several other examples are discussed in our companion paper.\\


%\noindent
%\textit{Israel-Stewart representation:} 
\section{Israel-Stewart representation}
In classical field theory, the existence of conservation laws is associated with the presence of a collection of currents $j_I^\mu$ (where $I$ is a new multi-index spanning the conserved quantities) with vanishing divergence $\partial_\mu j_I^\mu=0$ \cite{Peskin_book} (e.g., baryon number conservation). In linear fluid theories characterized by $\{\varphi^A,E^\mu,\sigma\}$ such as those considered here, such currents manifest themselves through the existence of a (constant) matrix $\mathcal{N}^A_I$ such that $\mathcal{N}^A_I \sigma_{AB}=0$.
%\begin{equation}\label{conserviamodibrutto}    \mathcal{N}^A_I \sigma_{AB}=0 \, .\end{equation}
In fact, if we contract the field equations of the theory, $E^\mu_{AB}\partial_\mu \varphi^B=-\sigma_{AB}\varphi^B$, with $\mathcal{N}^A_I$, the right-hand side vanishes, and we recover the equations $\partial_\mu j_I^\mu=0$, with
\begin{equation}\label{martina}
    j_I^\mu = \mathcal{N}^A_I E^\mu_{AB} \varphi^B \, .
\end{equation}
In the Supplementary Material, we prove the following useful result:
\begin{theorem}\label{theoLind}
Let $\{\varphi^A,E^\mu,\sigma \}$ be a linear theory for which all the hypotheses of Theorem \ref{theo} hold, and $\Xi_{[AB]}=0$. If the  conservation laws $\partial_\mu ( \mathcal{N}^A_I E^\mu_{AB} \varphi^B)=0$ are all independent, then there is a one-to-one change of variables $\varphi^A \rightarrow \{ \mu^I,\Pi^a \}$ such that $E^\mu$, $j_I^\mu$, and $\sigma$ take the form (all matrices below are constant)
\begin{equation}\label{decompose}
    E^\mu = \dfrac{1}{2} E^\mu_{IJ} \, \mu^I \mu^J + E^\mu_{Ib} \, \mu^I \Pi^b +\dfrac{1}{2} E^\mu_{ab} \, \Pi^a \Pi^b \, ,\qquad \qquad 
    j^\mu_I =  E^\mu_{IJ} \,  \mu^J + E^\mu_{Ib} \, \Pi^b  \, , \qquad \qquad
    \sigma = \sigma_{ab} \, \Pi^a \Pi^b \, ,  \\
\end{equation}
where $E^\mu_{IJ}$, $E^\mu_{ab}$, and $\sigma_{ab}$ are symmetric matrices. If $\mathcal{N}^A_I$ accounts for all conservation laws, then $\sigma_{ab}$ is invertible. 
\end{theorem}
% \begin{proof}
% We can regard the collection of fields $\varphi^A$ as a map $\varphi^A : \mathbb{R}^{1+3}\rightarrow \mathbb{R}^{\mathfrak{D}}$, for some positive integer $\mathfrak{D}$. For a fixed value of the index $I$, we may also view $\mathcal{N}^A_I$ as an element of the vector space $\mathbb{R}^{\mathfrak{D}}$. If the conservation laws are all independent, the vectors $\mathcal{N}^A_I$ (for all $I$) form a linearly independent set. Therefore, we can complete them to a basis of $\mathbb{R}^{\mathfrak{D}}$, through the introduction of some other (constant) vectors $\mathcal{N}^A_a$. It follows that the fields $\varphi^A$ can always be expressed in the form $\varphi^A=\mathcal{N}^A_I \mu^I+\mathcal{N}^A_a \Pi^a$, where the linear combination coefficients $\{\mu^I, \Pi^a \}$ are hydrodynamic fields of their own right (which are smooth if and only if $\varphi^A$ is smooth). Since the expansion of a vector on a basis is necessarily unique, the change of variables $\varphi^A \rightarrow \{\mu^I, \Pi^a \}$ is one-to-one. Plugging $\varphi^A=\mathcal{N}^A_I \mu^I+\mathcal{N}^A_a \Pi^a$ into equations \eqref{second} and \eqref{martina}, we obtain equation \eqref{decompose}, where the symmetry of $E^\mu_{IJ}$, $E^\mu_{ab}$, and $\sigma_{ab}$ follows from the symmetry of $E^\mu_{AB}$ and $\sigma_{AB}$. Finally, suppose that there is a vector $V^a \neq 0$ such that $V^a\sigma_{ab}=0$. Then, one can show that $V^a\mathcal{N}^A_a\sigma_{AB}=0$ (use $\sigma_{ab}=\sigma_{AB}\mathcal{N}^A_a\mathcal{N}^B_b$, and $\sigma_{AB}\mathcal{N}^B_I=0$), i.e. the matrix $\mathcal{N}^A_I$ is ``missing'' one conservation law, which can be accounted for by extending $\mathcal{N}^A_I$ to include $V^a\mathcal{N}^A_a$ as an additional (linearly independent) column. 
% \end{proof}
Theorem \ref{theoLind} tells us that one can always rearrange the EOM so that their mathematical structure ``resembles'' Israel-Stewart theory. In fact, if $\sigma_{ab}$ is invertible, with matrix inverse $\sigma^{ab}$, we can express the field equations \eqref{cinque} in terms of the variables $\{ \mu^I,\Pi^a \}$ as
\begin{equation}\label{lindblomrepr}    
    \partial_\mu (E^\mu_{IJ} \,  \mu^J + E^\mu_{Ib} \, \Pi^b)= 0 \, ,\qquad
       \sigma^{ab} E^\mu_{bc} \,\partial_\mu \Pi^c +  \,\Pi^a  = -  \sigma^{ab} E^\mu_{Jb} \,\partial_\mu \mu^J \, ,
\end{equation}
which we will refer to as the ``Israel-Stewart representation'' of the theory. The first set of equations in \eqref{lindblomrepr} is the set of all conservation laws. The second set gives relaxation-type equations \cite{Jou_Extended}, which describe dissipation. Stability requires that $E^0_{ab}$ and $\sigma_{ab}$ be positive definite so that the second equation in \eqref{lindblomrepr} contains both $\partial_t \Pi^a$ and $\Pi^a$, breaking time-reversal invariance. Thus, one can interpret $\Pi^a$ as ``dissipative fields'' \cite{LindblomRelaxation1996} (e.g., viscous stresses, diffusive currents, and reaction affinities \cite{PrigoginebookModernThermodynamics2014}). The fields $\mu^I$ are the usual ``dynamical fluid fields'' \cite{LindblomRelaxation1996} (e.g., temperature, chemical potential, and flow velocity). Indeed, if $\Pi^a=0$ then  $2E^\mu=\mu^I j_I^\mu$, meaning that $\mu^I$ may be interpreted as the ``chemical potential'' of the conserved density  $j_I^0$   \cite{GavassinoCasmir2022,Pathria2011}. 
Therefore, thermodynamically stable relativistic theories can always be expressed in this Israel-Stewart representation specified by a choice of $\{ \mu^I,\Pi^a,E^\mu,\sigma \}$.


%"distinguishing the set of equations of motion associated with conservation laws (i.e., energy, momentum, or other global charge densities) from those that are not, given the inability of the former to equilibrate on microscopic time scales."

Further insight can be obtained by realizing that, due to rotational invariance of the equilibrium state, we may further decompose $\mu^I$ and $\Pi^a$ into irreducible tensors of the $SO(3)$ rotation group. For example, if among the fields there is a 4-current $\delta n^\mu$, then $\delta n^0$ behaves as a scalar under rotations, and $\delta n^j$ ($j=1,2,3$) behaves as a 3-vector. Hence, we can assign to a given theory a list of integers $(\mathfrak{g}_0,\mathfrak{g}_1,\mathfrak{g}_2,...)$ that specifies the geometric character of its degrees of freedom,  with  $\mathfrak{g}_0$ being the number of scalars, $\mathfrak{g}_1$ the number of vectors, $\mathfrak{g}_2$ the number of symmetric traceless tensors with two indices, and so forth. One can repeat the same procedure for the conservation laws, including in the count only the fields $\mu^I$, which transform under $SO(3)$ as the respective conserved densities $j_I^0$. This produces a second list of integers,  $(\bar{\mathfrak{g}}_0,\bar{\mathfrak{g}}_1,\bar{\mathfrak{g}}_2,...)$, where $\bar{\mathfrak{g}}_n \leq \mathfrak{g}_n$. From \eqref{lindblomrepr}, we see that the theory is non-dissipative if and only if $\bar{\mathfrak{g}}_n =\mathfrak{g}_n$, $\forall \, n$, i.e., if there are as many conservation laws as degrees of freedom. Therefore, one can fully specify a given theory by saying that it is ``of class $(\mathfrak{g}_0,\mathfrak{g}_1,\mathfrak{g}_2,...)-(\bar{\mathfrak{g}}_0,\bar{\mathfrak{g}}_1,\bar{\mathfrak{g}}_2,...)$''. 
For example, a perfect fluid at finite chemical potential is of class $(2,1)-(2,1)$, and it does not dissipate, whereas a bulk-viscous fluid at zero chemical potential is of class $(2,1)-(1,1)$, and it dissipates (here, it is understood that $\mathfrak{g}_n=\bar{\mathfrak{g}}_n=0$, for $n \geq 2$). 
%To do this, we need to set up a systematic classification of fluid models into different ``types of matter'' (i.e. into different rheological models \cite{Krusser_2020}). According to this classification, two materials are considered fundamentally different if they have structurally different degrees of freedom, or different conservation laws.



%\noindent
%\textit{Universality classes}:  
\section{Universality classes}
%Our classification of fluid theories into universality classes considers only two ingredients: degrees of freedom and conservation laws. The idea is that, if two theories have analogous degrees of freedom and conservation laws, then they should be mathematically equivalent in the linear regime (apart from, maybe, postulating different values for some transport coefficients, see e.g. \cite{GavassinoGENERIC2022}). One can only show this on a case-by-case basis, by making clever use of Theorems \ref{theo}, \ref{theo2} and \ref{theoLind} (we will outline the procedure later). But if two linear theories present structurally different degrees of freedom or conservation laws, then they are fundamentally different, and describe ``types of matter'', which react differently to an imposed disturbance. To understand why we need both degrees of freedom and conservation laws in our classification, consider the two following questions:
%\begin{itemize}    \item What is the difference between a perfect fluid and a two-component plasma?
%  \item What is the difference between a two-component plasma and a superfluid?
%\end{itemize}
%In the first case, the difference is in the number of degrees of freedom, since we need two velocities to describe a two-component plasma, and only one to describe a perfect fluid. In the second case, however, both fluids have two velocities. What distinguishes them is the fact that a superfluid has an additional conservation law, namely the superfluid winding number \cite{AndreevMelnikovsky2004,Termo}, which is responsible for the long life of superfluid currents. Let us see how this example can be generalised to give rise to a systematic ``rheological classification'' of all types of substances.
%
%Given that our thermodynamic equilibrium state is homogeneous and isotropic,  the linearized hydrodynamic theory (as viewed in the equilibrium rest frame) must be invariant under the $SO(3)$ rotation group. As a consequence, the expressions for the information current $E^\mu$ and entropy production $\sigma$ cannot single out any preferred direction in space, imposing strong constraints on the structure of thermodynamically stable relativistic theories. For example, if our only degree of freedom is a 4-current $\delta n^\mu$, then at second-order in disturbances $E^0$ and $\sigma$ can only contain a term proportional to $(\delta n^0)^2$ and a term proportional to $\delta n^k \delta n_k$, while $E^j$ can only contain a term proportional to $\delta n^0 \delta n^j$. Therefore, one should classify degrees of freedom and conservation laws based on their geometric character. 
%Let $\{ \mu^I,\Pi^a,E^\mu,\sigma \}$ be a hydrodynamic theory expressed in the Israel-Stewart representation discussed above. Decompose its fields $\mu^I$ and $\Pi^a$ into irreducible tensors associated with the $SO(3)$ rotation group of the equilibrium state, and count how many tensors are present. For example, if among the fields there is a 4-current $\delta n^\mu$, then $\delta n^0$ should be counted as a scalar, and $\delta n^j$ ($j=1,2,3$) should be counted as a 3-vector. In this way, we can assign to the theory a list of integers $(\mathfrak{g}_0,\mathfrak{g}_1,\mathfrak{g}_2,...)$, where $\mathfrak{g}_0$ is the number of scalars, $\mathfrak{g}_1$ is the number of vectors, $\mathfrak{g}_2$ is the number of symmetric traceless tensors with two indices, and so on. The list $(\mathfrak{g}_0,\mathfrak{g}_1,\mathfrak{g}_2,...)$ classifies the degrees of freedom of the theory, since it counts both fields of ``$\mu^I$ type'' and fields of ``$\Pi^a$ type''. To classify the conservation laws, we can repeat the procedure above, including in the count only the fields $\mu^I$, which transform under $SO(3)$ like the respective conserved densities $j_I^0$. This produces a second list of integers,  $(\bar{\mathfrak{g}}_0,\bar{\mathfrak{g}}_1,\bar{\mathfrak{g}}_2,...)$, where $\bar{\mathfrak{g}}_n \leq \mathfrak{g}_n$. Then, we say that the theory is ``of class $(\mathfrak{g}_0,\mathfrak{g}_1,\mathfrak{g}_2,...)-(\bar{\mathfrak{g}}_0,\bar{\mathfrak{g}}_1,\bar{\mathfrak{g}}_2,...)$''. From \eqref{lindblomrepr}, we see that the theory is non-dissipative if and only if $\bar{\mathfrak{g}}_n =\mathfrak{g}_n$, $\forall \, n$, i.e. if there are as many conservation laws as degrees of freedom.
%For example, a perfect fluid at finite chemical potential is of class $(2,1)-(2,1)$, and it does not dissipate, whereas a bulk-viscous fluid at zero chemical potential is of class $(2,1)-(1,1)$, and it dissipates (here, it is understood that $\mathfrak{g}_n=\bar{\mathfrak{g}}_n=0$, for $n \geq 2$). 
%
%Furthermore, Theorems \ref{theo} and \ref{theo2} can be used to explain why all second order theories for viscous hydrodynamic must be indistinguishable in the linear regime. This will give to the reader a first taste of what we mean here by ``universality class''. In fact, one can show that the linearized Israel-Stewart theory in the Eckart frame at finite chemical potential \cite{Hishcock1983} is the most general linear theory that satisfies the hypotheses of Theorems \ref{theo} and \ref{theo2} and possesses the following properties: (i) it possesses 3 scalar degrees of freedom, 2 vector degrees of freedom, and 1 symmetric traceless tensor degree of freedom with two indices (this classification refers to the behavior of the fields $\varphi^A$ under rotations in the equilibrium rest frame, not under Lorentz transformations); (ii) particles, energy, and momentum are conserved; (iii) the stress-energy tensor is symmetric, and all the fields contribute explicitly to it (i.e. all $\varphi^A$ enter its constitutive relation). The proof of this statement is lengthy, and we will present it in our companion paper. However, we note that this result implies that divergence-type theories \cite{Geroch1995}, Carter's theory for viscosity \cite{carter1991}, BRSSS \cite{Baier2008}, DNMR \cite{Denicol2012Boltzmann} and the GENERIC theory \cite{Stricker2019} must all become indistinguishable from the Israel-Stewart theory near equilibrium. 
%
%The main goal of this Letter is to generalise the above result to different hydrodynamic systems (e.g., to fluid mixtures, superfluids, and magnetofluids). In order to do this, we need one last theorem.
Pick a class, i.e. fix the values of $\mathfrak{g}_n$ and $\bar{\mathfrak{g}}_n$. Working in the Israel-Stewart representation, give some names to the fields $\mu^I$ and $\Pi^a$ and construct the most general expressions for $E^\mu$ and $\sigma$, of the form \eqref{decompose}, compatible with rotational invariance. This produces the most general theory of the given class. By Theorem \ref{theo2}, any other theory belonging to the same class must be a particular case of this general theory (for some specific choice of parameters) since a mapping of the form \eqref{potato} is guaranteed to exist ($E^\mu$ and $\sigma$ being the most general). 
This general theory is usually very complicated, as it possesses a plethora of free coefficients. Luckily, Theorem \ref{theo2} comes to our aid: one can reabsorb many transport coefficients through changes of variables, as this does not modify the dynamics of the system. A useful type of field redefinition is the ``change of hydrodynamic frame'': $\mu^I = \tilde{\mu}^I +\mathcal{R}^I_c \tilde{\Pi}^c$ and $\Pi^a = \mathcal{R}^a_c \tilde{\Pi}^c$ ($\mathcal{R}^I_c$ and $\mathcal{R}^a_c$ being constant matrices), which preserves the structure \eqref{decompose},  mapping Israel-Stewart representations into  Israel-Stewart representations. The goal is to find a transformation that maps the general theory into an already existing theory (whose physical interpretation is known), which plays the role of a ``representative'' of the class. If this happens, Theorem \ref{theo2} guarantees that any linear theory belonging to the class is a particular realization of the representative and exhibits the same physical behavior.

We have applied this method to some selected (parity invariant \cite{SperanzaChiral2021}) classes. The representatives are reported in Table \ref{tableI} (we have fixed $\mathfrak{g}_n=\bar{\mathfrak{g}}_n=0$, for $n>2$). An interesting pattern can be recognized. In the absence of vector conservation laws, the currents do not have inertia, and they can only diffuse. If we include one vector conservation law, which plays the role of the linear momentum, only then will the system behave like an actual fluid. If we include more vector conservation laws, the fluid can sustain multiple non-diffusive relative flows, and it behaves like a superfluid. If we include a tensor conservation law, the system can conserve the ``memory'' of the deformations it experiences, and it becomes elastic. Combine two vector conservation laws with one tensor conservation law, and the result is a superfluid-elastic system, i.e., a supersolid.
\begin{table}[h!]
\centering
\begin{tabular}{ | c |c c c| l| }
\hline
$\mathfrak{g}_0  \, \mathfrak{g}_1 \, \mathfrak{g}_2 $&$\bar{\mathfrak{g}}_0$ & $\bar{\mathfrak{g}}_1$ & $\bar{\mathfrak{g}}_2$ & Representatives \\
 \hline
 \hline
  $a \, \, \, \, 0 \, \,  \,\, 0$ & $\leq a$ & 0 & 0 & Chemistry \\ 
 \hline
  $a \, \, \, \, 1 \, \,  \,\, 0$ & $\leq a$ & 1 & 0 & Fluid mixture \cite{noto_rel,BulkGavassino}; Models for bulk viscosity \cite{Causality_bulk} \\ 
 \hline
$a \, \, \, \, a \, \,  \,\, 0$ &  $\leq a$ & $ \leq a$ & 0 & Carter multifluids \cite{Carter_starting_point,GavassinoStabilityCarter2022} \\ 
 \hline
  \multirow{3}{3em}{ $1 \, \, \, \, 1 \, \,  \,\, 0$ }  & 0 & 0 & 0 & Diffusion of a non-conserved density \\ 
    & 1 & 0 & 0 & Cattaneo model of diffusion \cite{cattaneo1958,Jou_Extended} \\ 
      & 1 & 1 & 0 & Perfect fluid at $\mu=0$; Barotropic perfect fluid \\ 
 \hline
  \multirow{2}{3em}{ $2 \, \, \, \, 1 \, \,  \,\, 0$ }  & 1 & 1 & 0 & Bulk viscous fluid at $\mu=0$ \\ 
      & 2 & 1 & 0 & Perfect fluid  \\ 
\hline
  \multirow{4}{3em}{ $2 \, \, \, \, 2 \, \,  \,\, 0$ } & 2 & 0 & 0 & Coupled diffusion of two conserved densities \cite{Onsager1931} \\ 
  & 1 & 1 & 0 & Heat conductive bulk viscous fluid at $\mu=0$ \\ 
    & 2 & 1 & 0 & Heat conductive fluid at $\mu \neq 0$ \cite{OlsonRegularCarter1990} \\ 
      & 2 & 2 & 0 & Relativistic superfluid \cite{cool1995} \\ 
 \hline
   \multirow{2}{3em}{ $1 \, \, \, \, 1 \, \,  \,\, 1$ }  & 1 & 1 & 0 & Maxwell material \cite{BaggioliHolography,GavassinoGENERIC2022} at $\mu=0$ \\ 
      & 1 & 1 & 1 & Elastic material at $\mu=0$ or at $T=0$ \cite{Schumaker1983,landau7}   \\ 
 \hline
$1 \, \, \, \, 1 \, \,  \,\, 2$ &  1 & 1 & 0 & Burgers material \cite{Malek2018} at $\mu=0$; MIS$^*$ \cite{KeYin2022} \\ 
 \hline
\multirow{4}{3em}{$3 \, \, \, \, 2 \, \,  \,\, 1$} &  1 & 1 & 0 & Israel-Stewart theory in a ``general frame'' \cite{NoronhaGeneralFrame2021} at $\mu = 0$  \\ 
&  2 & 1 & 0 & Israel-Stewart theory \cite{Israel_Stewart_1979,Hishcock1983}  at $\mu \neq 0$ \\ 
&  3 & 1 & 1 & Elastic heat conducting material  \\ 
&  3 & 2 & 1 & Supersolid \cite{Andreev1969,Sears2010}; Inner crust of neutron stars \cite{CarterSamuelsson2006}  \\ 
 \hline
\end{tabular}
\caption{Universality classes near equilibrium.}
\label{tableI}
\end{table}
%\onecolumngrid
%\twocolumngrid
We provide the information currents and entropy production rates of the theories listed in Table \ref{tableI} in the Supplemental Material. In our companion paper, we use those expressions to discuss in detail the equivalence between the theories in the most relevant classes. We consider here
%However, if we include more vector conservation laws, the fluid can sustain multiple non-diffusive relative flows behaving as a superfluid. On the other hand, if  instead we include a tensor conservation law, the system can conserve the ``memory'' of the deformations it experiences, becoming an elastic material. Combine two vector conservation laws and one tensor conservation law and the result is a superfuid-elastic system, i.e. a supersolid.
%Relativistic supersolid hydrodynamics
a system of class $(1,1,1)-(1,1,1)$. Its fields are $\mu^I=\{\delta \mu, \delta u^k,\delta \Pi^{kl} \}$, which may be interpreted as the perturbations to the chemical potential, flow velocity, and shear stress tensor (which is symmetric and traceless). Since we have as many degrees of freedom as conservation laws (i.e., there are no fields of ``$\Pi^a$ type''), the entropy production rate vanishes. Thus, the most general theory is
\begin{equation}\label{infoelas}
    TE^0 =\dfrac{1}{2}  \dfrac{d n}{d\mu} (\delta \mu)^2 + \dfrac{1}{2} (\rho {+} P)\delta u^k \delta u_k + \dfrac{\delta \Pi^{kl}\delta \Pi_{kl}}{4G} \, ,\qquad
    TE^j = n\delta \mu \delta u^j + \delta \Pi^{jk} \delta u_k \, ,\qquad
    T\sigma = 0 \,.
\end{equation}
This information current contains all possible terms allowed by symmetry, making this theory a valid representative of its class. The transport coefficients have been given a physically-motivated name to ease their interpretation: $\rho {+} P$ may be interpreted as the enthalpy density, $G$ as the shear modulus, and $n$ as the background density. The prefactor of $\delta \Pi^{jk} \delta u_k$ could be set to unity by appropriately rescaling $\delta \Pi^{jk}$ with a field redefinition of the form $\delta \Pi^{jk} \rightarrow a \, \delta \Pi^{jk}$\footnote{The case in which the prefactor of $\delta \Pi^{jk} \delta u_k$ equals zero is trivial since it would produce the field equation $\partial_t \delta \Pi_{kl}=0$.}. If we compute the field equations from equation \eqref{cinque}, we obtain
\begin{equation}\label{dayofrain}
     \partial_t \delta n + \partial_j (n \delta u^j)=0 \, , \qquad
     (\rho {+} P)\partial_t \delta u_k + \partial_k (n \delta \mu) + \partial_j \delta \Pi^j_k =0 \, , \qquad
     \dfrac{\partial_t \delta \Pi_{kl}}{2G} + \braket{\partial_k \delta u_l}=0 \, , 
\end{equation}
where $\braket{...}$ extracts the symmetric traceless part. To obtain the last equation,  one needs to account for the constraints on $\delta \Pi_{kl}$ (see \cite{GavassinoCasmir2022} for a detailed discussion). Equations \eqref{dayofrain} describe the dynamics of an isotropic elastic material at zero temperature. The first equation is the continuity equation for particles, the second is the conservation of momentum, the third incorporates shear-stress dynamics in the Hookean approximation \cite{Schumaker1983}. Combining all three equations, and using $dP/d\rho = n^2 d\mu/(\rho {+} P)dn$, we obtain
\begin{equation}\label{anotherdayofsun}
    \partial^2_t \delta u_k -  \dfrac{G}{\rho {+} P} \partial^j \partial_j \delta u_k -\bigg[\dfrac{dP}{d\rho} + \dfrac{G}{3(\rho {+} P)} \bigg] \partial_k \partial_j \delta u^j =0 \, ,
\end{equation}
which is consistent with standard formulas from the theory of elasticity \cite{landau7}. This describes an elastic material that can sustain both longitudinal and transversal propagating waves with speeds $c_L=\sqrt{\dfrac{dP}{d\rho} + \dfrac{4G}{3(\rho {+}P)}}$ and $c_T= \sqrt{\dfrac{G}{\rho {+} P}}$, respectively \cite{landau7}.
%\begin{equation}
 %       c_L^2= \dfrac{dP}{d\rho} + \dfrac{4G}{3(\rho {+}P)} \, ,\qquad
 %       c_T^2= \dfrac{G}{\rho {+} P} \, ,
%\end{equation}
%which match the corresponding formulas from the theory of elasticity \cite{landau7}.

%Let us now comment on the conservation laws. Besides particles and momentum, the class $(1,1,1)-(1,1,1)$ displays a symmetric traceless tensor conservation law which (for elastic media) is emergent, and is analogous to the conservation of superflow in a superfluid \cite{Termo,AndreevMelnikovsky2004}. To understand its physical origin, consider the following. An elastic medium has an ``unstrained'' state, which is attained at equilibrium. A deviation from equilibrium corresponds to a displacement of the material elements from such unstrained state. Let $\xi_k$ be the corresponding displacement vector. In the linear regime, $\partial_t \xi_k=\delta u_k$ \cite{Schumaker1983}, and the third equation of \eqref{dayofrain} becomes
%$\partial_t (\delta \Pi_{kl} + 2G \! \braket{\partial_k \xi_l})=0$. Thus,  Hooke's law, $\delta \Pi_{kl}=-2G \! \braket{\partial_k \xi_l}$, is recovered when we integrate in time the additional tensor conservation law, and it manifests the ability of the system to retain (or, better, \textit{conserve}) the memory of its past deformations. 
%The associated conserved charge is the total deformation: $\mathbb{D}_{kl}={-}\int \!  \delta \Pi_{kl} d^3 x/2G = \oint \!  \braket{\xi_k d^2 S_l}$.


%\noindent
%\emph{Conclusions:} 
\section{Conclusions}
Our results establish, for the first time, connections between elasticity, viscosity, superfluidity, supersolidity, diffusion, and chemistry within a systematic, fully relativistic formalism. Using this paper, one can easily write down the (linearized) EOM of an arbitrary relativistic hydrodynamic system with very limited knowledge about its behavior: one only needs to know the relevant degrees of freedom and conservation laws \footnote{We note that, in general, knowledge about dispersion relations is not enough to specify the dynamics of a system uniquely. In fact, causality cannot solely be determined from dispersion relations, see \cite{Gavassino:2023mad} and our companion paper. On the other hand, within our approach, given the couple $\{E^\mu,\sigma\}$, the equations of motion are unique, and causal by construction. Furthermore, when one specifies the numbers $(\mathfrak{g}_0,\mathfrak{g}_1,\mathfrak{g}_2,...)-(\bar{\mathfrak{g}}_0,\bar{\mathfrak{g}}_1,\bar{\mathfrak{g}}_2,...)$, the transformation laws of the fields and possible constraints (such as $\Pi_{jk}=\Pi_{kj}$) are implemented from the start. Hence, our method, in general, provides more information about the system's dynamics than a pure spectral analysis.}. Then, the most general information current can be constructed, from which the EOM can be derived in full analogy with action principles. Our method's three main advantages are: (a) The resulting EOMs are always causal, stable, thermodynamically consistent, and uniquely solvable for smooth initial data. (b) Given the degrees of freedom and the conservation laws, the associated theory is unique. (c) The information current uniquely determines also the fluctuating generalization of the theory, giving rise to well-posed stochastic dynamics in terms of a path integral \cite{Mullins:2023tjg}.
 

%For a given collection of degrees of freedom, and a given set of conservation laws, what are the mechanical properties of the resulting fluid model near equilibrium that are compatible with relativity? This fundamental question is answered in this Letter through the combination of our Theorems \ref{theo}, \ref{theo2}, and \ref{theoLind} which show that the second law of thermodynamics and thermodynamic stability criteria are so strong that the very structure of the linearized fluid equations is fully determined by the choice of degrees of freedom and conservation laws. As a consequence, all thermodynamically consistent theories fall into novel types of universality classes near equilibrium. This marks the dawn of the ``relativistic rheology" era in hydrodynamics, where novel connections among a myriad of seemingly different physical systems are possible.  



%Our Theorems \ref{theo}, \ref{theo2}, and \ref{theoLind} show that the second law of thermodynamics and the stability criteria are so strong that the structure of the  linearized fluid equations is fully determined by the choice of degrees of freedom and conservation laws. As a consequence, all thermodynamically consistent theories fall into universality classes near equilibrium. This Letter sets the foundations of a new branch of relativistic hydrodynamics that we call ``relativistic rheology'', which provides systematic answers to questions such as: ``For a given collection of degrees of freedom, and a given set of conservation laws, what are the mechanical properties of the resulting fluid model near equilibrium?''


%This Letter sets the foundations of a new branch of relativistic hydrodynamics that may be called ``relativistic rheology'', as it aims to answer the following question: For a given collection of degrees of freedom, and a given set of conservation laws, what are the mechanical properties of the resulting fluid model near equilibrium? What makes a universal answer to this question possible are our Theorems \ref{theo}, \ref{theo2}, and \ref{theoLind}, which show that the second law of thermodynamics and the criteria for thermodynamic stability are so strong, that the structure of the linearised equations is fully determined by degrees of freedom and conservation laws, for all thermodynamically consistent theories to fall into universality classes near equilibrium. 


%Thus, following what is done in rheology \cite{Krusser_2020,Findley_book,Malkin_book}, one can construct a complete list of all possible universality classes and interpret each class just as a model for a specific ``type of material'', where different materials evolve in ways that are measurably different. Therefore, if two physicists aim to describe the same material, but the first is guided by, say, a thermodynamic principle, and the second is guided by, say, a variational principle, the resulting theories may be different, but (if everything is done correctly) they should fall into the same universality class, and should behave identically  close to equilibrium.

% \begin{equation}
%     \begin{split}
%         & \dfrac{\partial s}{\partial \mu}\bigg|_{T}  = \dfrac{\partial n}{\partial T}\bigg|_{\mu}   \, ,\\
%         & Ts+\mu n = \rho {+} P = \rho_S +\rho_N \, , \\
%         & {\mathcal{K}^{XY}= 
% \begin{bmatrix}
%   \dfrac{1}{s^2}\bigg(Ts-\mu n + \dfrac{\mu^2n^2}{\rho_S}\bigg) & \, \, \dfrac{\mu}{s} \bigg( 1-\dfrac{\mu n}{\rho_S} \bigg)  \\
%   \dfrac{\mu}{s} \bigg( 1-\dfrac{\mu n}{\rho_S} \bigg) &  \dfrac{\mu^2}{\rho_S} \\
% \end{bmatrix} ,} \\
%     \end{split}
% \end{equation}

%\noindent
%\emph{Acknowledgements:}
\section{Acknowledgements}
M.M.D. is partially supported by NSF grant DMS-2107701, a Vanderbilt's Seeding Success Grant, 
a Chancellor's Faculty Fellowship, and DOE grant DE-SC0024711.
L.G. is partially supported by a Vanderbilt's Seeding Success Grant. J.N. is partially supported by the U.S. Department of Energy, Office of Science, Office for Nuclear Physics
under Award No. DE-SC0023861.

\bibliography{Biblio}


\onecolumngrid
\newpage
\begin{center}
{\bf \large SUPPLEMENTARY MATERIAL}
\end{center}

\setcounter{equation}{0}
\setcounter{figure}{0}
\setcounter{table}{0}
\setcounter{page}{1}
\renewcommand{\theequation}{S\arabic{equation}}
\renewcommand{\thefigure}{S\arabic{figure}}
%\renewcommand{\bibnumfmt}[1]{[S#1]}
%\renewcommand{\citenumfont}[1]{S#1}



\maketitle


\section{Origin and properties of the information current}\label{sectionI}

Here, we briefly review how the information current is explicitly computed for an assigned (fully non-linear) theory. The foundations of the method where laid out in \cite{GavassinoGibbs2021} and refined in subsequent works \cite{GavassinoCausality2021,
GavassinoStabilityCarter2022,GavassinoGENERIC2022}. 

The equilibrium statistical operator of a fluid in contact with a large environment is $\hat{\rho}\propto \exp(\alpha_\star^I \hat{Q}_I)$, where $\hat{Q}_I$ are the Noether charges (energy, momentum, electric charge,...) of the underlying microscopic theory, and $\alpha_\star^I$ are some constants which are uniquely determined by the thermodynamic state of the environment. The operator $\hat{\rho}$ defines the grand-canonical ensemble of a system with arbitrary conservation laws \cite{Hakim2011}. Let $\psi^A$ be the hydrodynamic fields of the (fully non-linear) hydrodynamic theory, and consider a specific field configuration $\psi^A(\Sigma)$  on a given Cauchy surface $\Sigma$. This configuration constitutes a possible choice of initial data for the system, and defines a macroscopic state, with many possible microscopic realizations. Let $\hat{\mathcal{P}}[\psi^A(\Sigma)]$ be the projector onto the Hilbert subspace of all the possible microscopic realizations of $\psi^A(\Sigma)$. The entropy of $\psi^A(\Sigma)$ is given by Boltzmann's rule: $\exp S[\psi^A(\Sigma)]=\text{Tr} \hat{\mathcal{P}}[\psi^A(\Sigma)] $. Therefore, the equilibrium probability of observing the hydrodynamic state $\psi^A(\Sigma)$ is
\begin{equation}
\mathcal{P} \propto \text{Tr}\big(\hat{\mathcal{P}}[\psi^A(\Sigma)]\exp(\alpha_\star^I \hat{Q}_I) \big) \approx \text{Tr}\big(  \hat{\mathcal{P}}[\psi^A(\Sigma)] \big) \exp(\alpha_\star^I Q_I[\psi^A(\Sigma)]) = \exp(S[\psi^A(\Sigma)]+\alpha_\star^I Q_I[\psi^A(\Sigma)]) \, .
\end{equation}
Clearly, the equilibrium macrostate is the most probable macrostate. We can therefore conclude that the equilibrium state of the hydrodynamic theory must be the state which maximizes (for unconstrained variations) the functional
\begin{equation}
\Phi[\psi^A]= S[\psi^A] + \alpha_\star^I Q_I[\psi^A] \, ,
\end{equation} 
for fixed values of all $\alpha_\star^I$ (which depend on the environment). Then, assuming that the theory has an entropy current $s^\mu(\psi^A)$, we need to maximize the flux-integral (we adopt the ``standard orientation'' \cite{MTW_book}: $d\Sigma_0 >0$)
\begin{equation}\label{phiuzzo}
\Phi= \int_\Sigma (s^\mu +\alpha_\star^I J^\mu_I)d\Sigma_\mu  \, ,
\end{equation}
where $J^\mu_I(\psi^A)$ are the conserved currents of the theory. In order to actually compute the maximum of \eqref{phiuzzo}, we introduce a smooth one-parameter family, $\psi^A(\epsilon)$, of solutions of the field equations such that $\epsilon=0$ is the equilibrium state. Then, we impose $\dot{\Phi}(\epsilon=0)=0$ for any possible one-parameter family of this kind (we adopt the notation $\dot{f}:=df/d\epsilon$). This immediately leads to the covariant Gibbs relation \cite{Israel_2009_inbook},
\begin{equation}\label{Gibbsusss}
\dfrac{\partial s^\mu}{\partial \psi^A} = -\alpha_\star^I \dfrac{\partial J^\mu_I}{\partial \psi^A} \, ,
\end{equation}
which fixes the equilibrium state uniquely. Then, if we define $E:=\Phi(0)-\Phi(\epsilon)$, the condition that $\Phi$ should be maximized in equilibrium is equivalent to the requirement that $E$ should be strictly positive for non-vanishing perturbations. Performing a second-order Taylor expansion in $\epsilon$, and recalling equation \eqref{Gibbsusss}, we obtain\footnote{Note that the terms proportional to $\ddot{\psi}^A$ in $E$ cancel out as a consequence of \eqref{Gibbsusss}. }
\begin{equation}
E = \int_\Sigma -\dfrac{1}{2}\bigg[\dfrac{\partial^2 s^\mu}{\partial \psi^A \partial \psi^B} + \alpha_\star^I\dfrac{\partial^2 J^\mu_I}{\partial \psi^A \partial \psi^B}  \bigg] \varphi^A \varphi^B d\Sigma_\mu +\mathcal{O}(\epsilon^3) \geq 0  \, .
\end{equation}
The quantity inside the integral is $E^\mu$, and we are adopting the notation $\varphi^A:= \psi^A(\epsilon)-\psi^A(0)=\dot{\psi}^A(0) \epsilon+\mathcal{O}(\epsilon^2)$. If we require $E$ to be positive-definite for all choices of $\Sigma$, we find that $E^\mu$ is future-directed timelike. This automatically enforces hydrodynamic stability (in the linear regime), if the theory is consistent with the second law of thermodynamics. To see this, consider that, if no interaction occurs between the fluid and the environment, then $\Delta S \geq 0$ and $\Delta Q_I=0$ along any hydrodynamic process. Therefore, $\Delta \Phi=\Delta S \geq 0$ (recall that $\alpha_\star^I$ are constants). It follows that $\Phi(\epsilon)$ can only increase, or stay constant. Considering that $\Phi(0)$ is the equilibrium value of $\Phi$ (which is a constant), we can conclude that $E=\Phi(0)-\Phi(\epsilon)$ is non-increasing in time. In summary, $E$ plays the role of a positive-definite square-integral norm of the perturbations, which is bounded from below by zero and from above by its initial value. This guarantees (Lyapunov) stability of the equilibrium and global uniqueness of the solution for square-integrable initial data, in the linear regime \cite{Wald}. With a simple Gauss-theorem argument, one can also prove causality \cite{GavassinoCausality2021}.










\newpage

% \section{Proof of Theorem 1}


% \begin{proof}
% The system of partial differential equations on $\mathbb{R}^{1+3}$,
% \begin{equation}\label{prima}
% M^\mu_{AB}\partial_\mu \varphi^B = -\Xi_{AB} \varphi^B \, ,
% \end{equation}
% where $M^\mu_{AB}$ and $\Xi_{AB}$ are constant
% matrices, and $M^0_{AB}$ is invertible, admits smooth solutions of the form
% \begin{equation}\label{thevictory}
% \begin{split}
%     \varphi^A(t,x^j)={}&  \big(e^{-(M^0)^{-1} \Xi t} \big)\indices{^A_B} \, (Z^B-W^B) \\
%     +{}& \big(e^{-(M^0)^{-1} (\Xi+M^j a_j) t} \big)\indices{^A_B} \, W^Be^{a_j x^j} \, , \\
%     \end{split}
% \end{equation}
% for any $Z^A,W^A,a_j=\text{const}$. Thus, $E^\mu_{AB}Z^AZ^B= 2E^\mu(0)$ must be timelike-future directed for any $Z^A \neq 0$. In particular, $E^0_{AB}$ must be positive-definite (and, thus, invertible). Now, we always have the freedom to redefine the matrices $M^\mu_{AB}$ and $\Xi_{AB}$ by contracting both sides of \eqref{prima} with an invertible matrix $\mathcal{N}^A_C$. Let us then ``fix'' the matrix $M^0_{AB}$ to coincide with $E^0_{AB}$ (this is possible because both are invertible). If we contract \eqref{prima} with $\varphi^A$, and we plug 
% \begin{equation}\label{second}
% \begin{split}
%  E^\mu={}& \dfrac{1}{2} E^\mu_{AB} \varphi^A \varphi^B \, ,  \\
%  \sigma ={}&  \sigma_{AB} \varphi^A \varphi^B \, , \\
% \end{split}
% \end{equation}
% into \begin{equation}\label{Emu}
% \partial_\mu E^\mu = -\sigma \leq 0 \, ,
% \end{equation} we obtain the two equations below:
% \begin{equation}\label{gringo}
% \begin{split}
% & \varphi^A M^\mu_{AB}\partial_\mu \varphi^B + \varphi^A \Xi_{AB} \varphi^B =0 \, , \\
% &  \varphi^A E^\mu_{AB} \partial_\mu \varphi^B + \varphi^A \sigma_{AB} \varphi^B =0 \, . \\
% \end{split}
% \end{equation}
% Both are respected along \textit{all} solutions of \eqref{prima}. If we subtract the second equation of \eqref{gringo} to the first, the terms $\varphi^A M^0_{AB}\partial_t \varphi^B$ cancel out (we have fixed $M^0_{AB}=E^0_{AB}$). Evaluating the result along \eqref{thevictory}, at $x^\mu=0$, we obtain
% \begin{equation}\label{MjjjEjjj}
% Z^A(M^j_{AB}-E^j_{AB})a_j W^B +Z^A (\Xi_{AB}- \sigma_{AB} ) Z^B =0 \, .
% \end{equation}
% Since this must be true for any choice of $Z^A$, $W^A$, and $a_j$, we obtain $M^j_{AB}=E^j_{AB}$ and $\Xi_{(AB)}= \sigma_{AB}$. Hence, we have recovered \begin{equation}\label{cinque}
% E^\mu_{AB} \partial_\mu \varphi^B = -\sigma_{AB} \varphi^B - \Xi_{[AB]}\varphi^B \,.  
% \end{equation}
% On the other hand, the matrices $E^\mu_{AB}$ are symmetric, and $E^0_{AB}$ is positive definite. Thus, the system \eqref{cinque} is symmetric hyperbolic \footnote{Note that, if in \eqref{second} we did not define $E^\mu_{AB}$ to be symmetric, in the second line of \eqref{gringo} there would be $E^\mu_{(AB)}$ in place of $E^\mu_{AB}$, and we would eventually get $M^\mu_{AB}=E^\mu_{(AB)}$.} and, since $E^\mu_{AB}Z^AZ^B$ is future-directed timelike for any $Z^A \neq 0$, it is also causal \cite{Geroch_Lindblom_1991_causal}.
% \end{proof}

\section{Understanding Theorem 1} 

Theorem 1 is quite a surprising result, as it tells us that the very existence of the information current fully constrains the mathematical structure of the field equations. How is it possible? Here, we present an alternative proof of Theorem 1, which should help the reader understand what is really going on. We also provide an explicit example.


\subsection{Proof from linearity}

It is possible to prove Theorem 1 without ever manipulating the field equations ``$M^\mu_{AB}\partial_\mu \varphi^B = -\Xi_{AB}\varphi^B$'' explicitly, but only relying on their linearity. Let us see how this works. Consider two arbitrary solutions, $\varphi^A$ and $\tilde{\varphi}^A$, of the field equations. Then, since $\partial_\mu E^\mu =-\sigma$ holds along all solutions of the field equations, we must have\footnote{Note that, if we did not define $E^\mu_{AB}$ to be symmetric, in equation \eqref{pippo} there would be $E^\mu_{(AB)}$ in place of $E^\mu_{AB}$.} (recalling that $2E^\mu=E^\mu_{AB}\varphi^A \varphi^B$ and $\sigma=\sigma_{AB}\varphi^A \varphi^B$)
\begin{equation}\label{pippo}
\begin{split}
& \varphi^A E^\mu_{AB}\partial_\mu \varphi^B +\varphi^A \sigma_{AB}\varphi^B =0 \, , \\
& \tilde{\varphi}^A E^\mu_{AB}\partial_\mu \tilde{\varphi}^B +\tilde{\varphi}^A \sigma_{AB}\tilde{\varphi}^B =0 \, . \\
\end{split}
\end{equation}
Furthermore, since the equations are linear, also $\varphi^A+  \tilde{\varphi}^A$ must be a solution of the field equations, and $\partial_\mu E^\mu =-\sigma$ (which is obeyed by \textit{all} solutions) must hold also along $\varphi^A+  \tilde{\varphi}^A$. Therefore:
\begin{equation}\label{varuzzo}
(\varphi^A+ \tilde{\varphi}^A) E^\mu_{AB}\partial_\mu (\varphi^B+ \tilde{\varphi}^B) +(\varphi^A+ \tilde{\varphi}^A) \sigma_{AB}(\varphi^B+ \tilde{\varphi}^B) =0 \, .
\end{equation}
Expanding all the products, and using \eqref{pippo} to cancel some terms, we obtain
\begin{equation}\label{iotiuso}
\tilde{\varphi}^A E^\mu_{AB}\partial_\mu \varphi^B +\varphi^A E^\mu_{AB}\partial_\mu \tilde{\varphi}^B +2 \, \tilde{\varphi}^A \sigma_{AB}\varphi^B  =0 \, .
\end{equation}
Now we immediately see what is happening: since $E^\mu$ is quadratic in the fields, and the second law $\partial_\mu E^\mu =-\sigma$ must hold for all linear combinations of solutions, there are some ``cross-conditions'' on all couples of solutions $\{\varphi^A, \tilde{\varphi}^A \}$, and this induces very strong constraints on the structure of the field equations. For example, if at some point $p$ we have $\varphi^A(p)=0$, then equation \eqref{iotiuso} implies that $E^\mu_{AB}\partial_\mu \varphi^B=0$ at $p$. We can already see that the matrices $E^\mu_{AB}$ must be somehow related to the matrices $M^\mu_{AB}$ defining the principal part of the field equations. Our goal now is to show that we can fully ``reconstruct'' the field equations of the theory just using \eqref{iotiuso}. The idea is the following.

Suppose that the multi-index $A$ runs from $1$ to $\mathfrak{D}$. Pick an arbitrary spacetime event $p$ and construct a collection of $\mathfrak{D}$ local solutions\footnote{Note that the existence of such local solutions is guaranteed by the Cauchy-Kovalevskaya (CK) theorem \cite{Rauch_book,CourantHilbert2_book}. In fact, the invertibility of $M^0_{AB}$ guarantees that the equations can be recast in the normal form, and linear equations with constant coefficients obviously have analytic coefficients, so that the CK theorem applies. The exponential solutions that we present in the main text are indeed an example of analytic solutions. It is also easy to show that, if we release the assumption that $M^0_{AB}$ is invertible, the proof no longer works. For example, take the case $M^\mu_{AB}=0$, $\Xi_{AB}=\delta_{AB}$. Then the only solution is $\varphi^A=0$, and there exists no solution $\tilde{\varphi}^A$ such that $\tilde{\varphi}^A(p)=\delta^A_C$.}, $\{\tilde{\varphi}^A_{(C)} \}_{C=1}^\mathfrak{D}$, of the field equations, such that $\tilde{\varphi}^A_{(C)}(p)=\delta^A_C$. Then, since $\tilde{\varphi}^A_{(C)}$ are solutions of the field equations, equation \eqref{iotiuso} must hold also replacing $\tilde{\varphi}^A$ with $\tilde{\varphi}^A_{(C)}$. This produces a collection of $\mathfrak{D}$ conditions at $p$:
\begin{equation}
E^\mu_{CB}\partial_\mu \varphi^B +\varphi^A E^\mu_{AB}\partial_\mu \tilde{\varphi}^B_{(C)} +2 \, \sigma_{CB}\varphi^B  =0 \spc (\text{at }p).
\end{equation}
Defining $\mathcal{G}_{CA}:= [E^\mu_{AB}\partial_\mu \tilde{\varphi}^B_{(C)}]_p$ (the subscript $p$ means ``evaluated at $p$''), and isolating $E^\mu_{CB}\partial_\mu \varphi^B$, we obtain
\begin{equation}\label{almost}
E^\mu_{CB}\partial_\mu \varphi^B =-(2 \, \sigma_{CB}+\mathcal{G}_{CB})\varphi^B   \spc (\text{at }p).
\end{equation}
On the other hand, we know that the linearized theory is invariant under spacetime translations (because the background state is homogeneous). Hence, we can repeat the same procedure above at any other spacetime event $q$, operating a translation $\tilde{\varphi}^A_{(C)}(x)\rightarrow \tilde{\varphi}^A_{(C)}(x-q+p)$, so that now they reduce to $\delta^A_C$ at $q$. This implies that the condition \eqref{almost} must hold at every spacetime point, with the same $\mathcal{G}_{CB}$. Furthermore, if we plug \eqref{almost} into the first line of \eqref{pippo}, we find that the matrix $\sigma_{CB}+\mathcal{G}_{CB}$ must be anti-symmetric. Hence, we call it $\Xi_{[CB]}$, and \eqref{almost} finally becomes
\begin{equation}\label{finally}
E^\mu_{CB}\partial_\mu \varphi^B =-\sigma_{CB} \varphi^B -\Xi_{[CB]}\varphi^B \spc (\text{at any event }q) \, .
\end{equation}
Since $E^\mu_{CB}$ are all symmetric, and $E^\mu_{CB}\varphi^C \varphi^B=2E^\mu$ is future directed timelike $\forall \, \varphi^B\neq 0$, this is a  symmetric hyperbolic and  causal \cite{Geroch_Lindblom_1991_causal} system of $\mathfrak{D}$ partial differential equations in $\mathfrak{D}$ variables. Holmgren's uniqueness theorem \cite{Rauch_book} guarantees that such system uniquely fixes the fields for given initial data. Hence, it must necessarily contain as much information as the original field equations (which were not even invoked explicitly!). This completes our alternative proof of Theorem 1.

\subsection{Example: Deriving Maxwell's equations from  energy conservation}


Suppose that one has forgotten the two ``dynamical'' Maxwell equations in vacuum, namely $\partial_t \textbf{E}=\nabla \times \textbf{B}$, and $\partial_t \textbf{B}=-\nabla \times \textbf{E}$. However, one remembers that they are linear, and that they obey a conservation law of the form $\partial_t \rho + \nabla \cdot \textbf{S}=0$, where $\rho$ is the electromagnetic energy density and $\textbf{S}$ is the Poynting vector:
\begin{equation}
\rho = \dfrac{|\textbf{E}|^2 + |\textbf{B}|^2}{2} \, , \spc \textbf{S}= \textbf{E} \times \textbf{B} \, .
\end{equation}
Can one reconstruct the original equations? Let's see. We consider two arbitrary solutions of the (unknown) Maxwell's equations. We call them $\{\textbf{E},\textbf{B} \}$ and $\{ \tilde{\textbf{E}},\tilde{\textbf{B}} \}$. Then, the conservation law $\partial_t \rho + \nabla \cdot \textbf{S}=0$ implies
\begin{equation}\label{eormf}
\begin{split}
& \textbf{E} \cdot \partial_t \textbf{E} + \textbf{B} \cdot \partial_t \textbf{B} + \textbf{B} \cdot (\nabla \times \textbf{E})- \textbf{E} \cdot (\nabla \times \textbf{B}) =0 \, , \\
& \tilde{\textbf{E}} \cdot \partial_t \tilde{\textbf{E}} + \tilde{\textbf{B}} \cdot \partial_t \tilde{\textbf{B}} + \tilde{\textbf{B}} \cdot (\nabla \times \tilde{\textbf{E}})- \tilde{\textbf{E}} \cdot (\nabla \times \tilde{\textbf{B}}) =0 \, .\\
\end{split}
\end{equation}
This does not fix the evolution of $\textbf{E}$ and $\textbf{B}$ completely. However, let us now consider the sum $\{\textbf{E}+\tilde{\textbf{E}},\textbf{B}+\tilde{\textbf{B}} \}$. Since we know that Maxwell's equations in vacuum are linear, we know that this is also a solution. Hence, the conservation law $\partial_t \rho + \nabla \cdot \textbf{S}=0$ must hold also for the total energy density and flux:
\begin{equation}
\rho = \dfrac{|\textbf{E}+\tilde{\textbf{E}}|^2 + |\textbf{B}+\tilde{\textbf{B}}|^2}{2} \, , \spc \textbf{S}= (\textbf{E}+\tilde{\textbf{E}}) \times (\textbf{B}+\tilde{\textbf{B}}) \, .
\end{equation}
This produces the cross condition below:
\begin{equation}\label{gbuf}
\tilde{\textbf{E}} \cdot (\partial_t \textbf{E}-\nabla \times \textbf{B}) + \tilde{\textbf{B}} \cdot (\partial_t \textbf{B}+\nabla \times \textbf{E} ) +\textbf{E} \cdot (\partial_t \tilde{\textbf{E}}-\nabla \times \tilde{\textbf{B}}) + \textbf{B} \cdot (\partial_t \tilde{\textbf{B}}+\nabla \times \tilde{\textbf{E}} ) =0 \, .
\end{equation}
Now, fixed some event $p$, the values of $\tilde{\textbf{E}}$ and $\tilde{\textbf{B}}$ at $p$ may be arbitrary. Thus, if we want equation \eqref{gbuf} to hold for any choice of $\{ \tilde{\textbf{E}},\tilde{\textbf{B}} \}$, we must require that $\partial_t \textbf{E}-\nabla \times \textbf{B}$ and $\partial_t \textbf{B}+\nabla \times \textbf{E}$  be linear combinations of $\textbf{E}$ and $\textbf{B}$. Then, it is easy to see that the most general (rotationally-invariant) field equations that are consistent with \eqref{eormf} and \eqref{gbuf} are
\begin{equation}
\partial_t \textbf{E}-\nabla \times \textbf{B}= \mathcal{G}\, \textbf{B}  \, , \spc \partial_t \textbf{B}+\nabla \times \textbf{E}= -\mathcal{G} \, \textbf{E} \, ,
\end{equation}
where $\mathcal{G}$ is some constant. This a symmetric-hyperbolic system, in agreement with our theorem. As we can see, we could fully reconstruct Maxwell's equations from the condition $\partial_t \rho + \nabla \cdot \textbf{S}=0$, except for the antisymmetric coupling $\mathcal{G}$ (which corresponds to the matrix $\Xi_{[AB]}$, in our theorem). Nature has chosen $\mathcal{G}=0$.

\newpage

\section{Proof of Theorem 3}


\textbf{Theorem 3}. 
\textit{Let $\{\varphi^A,E^\mu,\sigma \}$ be a linear theory for which all the hypotheses of Theorem 1 hold, and $\Xi_{[AB]}=0$. If the  conservation laws $\partial_\mu ( \mathcal{N}^A_I E^\mu_{AB} \varphi^B)=0$ are all independent, then there is a one-to-one change of variables $\varphi^A \rightarrow \{ \mu^I,\Pi^a \}$ such that $E^\mu$, $j_I^\mu$, and $\sigma$ take the form (all matrices below are constant)}
\begin{equation}\label{decompose}
    E^\mu = \dfrac{1}{2} E^\mu_{IJ} \, \mu^I \mu^J + E^\mu_{Ib} \, \mu^I \Pi^b +\dfrac{1}{2} E^\mu_{ab} \, \Pi^a \Pi^b \, ,\qquad \qquad 
    j^\mu_I =  E^\mu_{IJ} \,  \mu^J + E^\mu_{Ib} \, \Pi^b  \, , \qquad \qquad
    \sigma = \sigma_{ab} \, \Pi^a \Pi^b \, ,  \\
\end{equation}
\textit{where $E^\mu_{IJ}$, $E^\mu_{ab}$, and $\sigma_{ab}$ are symmetric matrices. If $\mathcal{N}^A_I$ accounts for all conservation laws, then $\sigma_{ab}$ is invertible.} 
\begin{proof}
We can regard the collection of fields $\varphi^A$ as a map $\varphi^A : \mathbb{R}^{1+3}\rightarrow \mathbb{R}^{\mathfrak{D}}$, for some positive integer $\mathfrak{D}$. For a fixed value of the index $I$, we may also view $\mathcal{N}^A_I$ as an element of the vector space $\mathbb{R}^{\mathfrak{D}}$. If the conservation laws are all independent, the vectors $\mathcal{N}^A_I$ (for all $I$) form a linearly independent set. Therefore, we can complete them to a basis of $\mathbb{R}^{\mathfrak{D}}$, through the introduction of some other (constant) vectors $\mathcal{N}^A_a$. It follows that the fields $\varphi^A$ can always be expressed in the form $\varphi^A=\mathcal{N}^A_I \mu^I+\mathcal{N}^A_a \Pi^a$, where the linear combination coefficients $\{\mu^I, \Pi^a \}$ are hydrodynamic fields of their own right (which are smooth if and only if $\varphi^A$ is smooth). Since the expansion of a vector on a basis is necessarily unique, the change of variables $\varphi^A \rightarrow \{\mu^I, \Pi^a \}$ is one-to-one. Plugging $\varphi^A=\mathcal{N}^A_I \mu^I+\mathcal{N}^A_a \Pi^a$ into the constitutive relations
\begin{equation}
     E^\mu= \dfrac{1}{2} E^\mu_{AB} \varphi^A \varphi^B\,, \, \qquad  j_I^\mu = \mathcal{N}^A_I E^\mu_{AB} \varphi^B \, , \qquad 
 \sigma =  \sigma_{AB} \varphi^A \varphi^B \,, 
\end{equation}
we obtain equation \eqref{decompose}, where the symmetry of $E^\mu_{IJ}$, $E^\mu_{ab}$, and $\sigma_{ab}$ follows from the symmetry of $E^\mu_{AB}$ and $\sigma_{AB}$. Finally, suppose that there is a vector $V^a \neq 0$ such that $V^a\sigma_{ab}=0$. Then, one can show that $V^a\mathcal{N}^A_a\sigma_{AB}=0$ (use $\sigma_{ab}=\sigma_{AB}\mathcal{N}^A_a\mathcal{N}^B_b$, and $\sigma_{AB}\mathcal{N}^B_I=0$), i.e. the matrix $\mathcal{N}^A_I$ is ``missing'' one conservation law, which can be accounted for by extending $\mathcal{N}^A_I$ to include $V^a\mathcal{N}^A_a$ as an additional (linearly independent) column. 
\end{proof}

\newpage

\section{Atlas of Information Currents and entropy production rates}

Below, we provide the information currents and entropy production rates of the theories listed in the table present in the main text.

\noindent \textit{How to interpret this atlas -} Both $E^\mu$ and $\sigma$ have been rescaled by a constant factor $T$ (background temperature) for convenience. The fields $\varphi^A$ describe small displacements from equilibrium, which are marked with a variation symbol ``$\delta$'', e.g. $\delta T$, or $\delta u^k$.  The coefficients in front of them (i.e., all the quantities without ``$\delta$'') are background quantities, which must be treated as constants. In some information currents, we have introduced a ``chemical'' index $X,Y \in \{1,...,a\}$. Einstein's summation convention applies to all types of indices, including $X$ and $Y$. Chemical matrices, such as $P_{XY}$ and $\mathcal{K}^{XY}$, are all symmetric. 


\noindent \textit{How to use this atlas -} With the aid of Theorem 1, one can fully reconstruct the linear field equations of each theory directly from the formulas of $TE^\mu$ and $T\sigma$. In fact, the system $E^\mu_{AB}\partial_\mu \varphi^A = -\sigma_{AB}\varphi^B$ is formally equivalent to \cite{GavassinoNonHydro2022}
\begin{equation}\label{systemreads}
    \partial_\mu \dfrac{\partial (TE^\mu)}{\partial \varphi^A} = -\dfrac{1}{2} \dfrac{\partial (T\sigma)}{\partial \varphi^A}.
\end{equation}
For example, take the diffusion of a non-conserved density, class $(1,1,0)-(0,0,0)$. Plugging \eqref{diffusuz} into \eqref{systemreads}, we obtain
\begin{equation}\label{systemreads33}
\begin{split}
    \partial_\mu \dfrac{\partial (TE^\mu)}{\partial (\delta \mu)} ={}& -\dfrac{1}{2} \dfrac{\partial (T\sigma)}{\partial (\delta \mu)} \spc \Longrightarrow \spc \partial_t (A \delta \mu) + \partial_j \delta j^j =-\Xi \delta \mu \\
    \partial_\mu \dfrac{\partial (TE^\mu)}{\partial (\delta j^k)} ={}& -\dfrac{1}{2} \dfrac{\partial (T\sigma)}{\partial (\delta j^k)} \spc \Longrightarrow \spc \partial_t (\mathcal{K} \delta j_k) + \partial_k \delta \mu =-\mathcal{R} \delta j_k\\
    \end{split}
\end{equation}
As one can see, the system \eqref{systemreads33} is symmetric by construction. Once the field equations have been obtained, it is possible to assign an interpretation to the fields and the transport coefficients based on their role in the dynamics. For example, in the system \eqref{systemreads33}, we can interpret $A\delta \mu$ as the density of a non-conserved charge, $\delta \mu$ as its chemical potential (which should vanish at equilibrium), $\delta j^k$ as its flux, and $\Xi$ as its production rate. The second equation of \eqref{systemreads33} is a Cattaneo-type relaxation equation for the flux $\delta j^k$, with relaxation time $\mathcal{K}/\mathcal{R}$ and diffusivity coefficient $\mathcal{R}^{-1}$.



\subsection*{Chemistry: ($a,0,0$)-($\leq a, 0,0$)}

\begin{equation}
\begin{split}
TE^0 ={}& \dfrac{1}{2} P_{XY} \, \delta \mu^X \delta \mu^Y \, ,  \\
TE^j ={}& 0 \, , \\
T\sigma ={}& \Xi_{XY} \, \delta \mu^X \delta \mu^Y \, , \\
\end{split}
\end{equation}







\subsection*{Fluid mixtures: ($a,1,0$)-($\leq a, 1,0$)}

\begin{equation}
\begin{split}
TE^0 ={}& \dfrac{1}{2} P_{XY} \, \delta \mu^X \delta \mu^Y + \dfrac{1}{2}(\rho {+} P)\delta u_k \delta u^k \, ,\\
TE^j ={}& n_X \delta \mu^X \delta u^j \, , \\
T\sigma ={}& \Xi_{XY} \, \delta \mu^X \delta \mu^Y \, ,\\
\end{split}
\end{equation}



\subsection*{Carter multifluids: ($a,a,0$)-($\leq a, \leq a,0$)}

\begin{equation}
\begin{split}
TE^0 ={}& \dfrac{1}{2} P_{XY} \, \delta \mu^X \delta \mu^Y + \dfrac{1}{2}\mathcal{K}^{XY}\delta j_{Xk} \delta  j_Y^k \\
TE^j ={}&  \delta \mu^X \delta j_X^j \\
T\sigma ={}& \Xi_{XY} \, \delta \mu^X \delta \mu^Y + \mathcal{R}^{XY} \delta j_{Xk} \delta  j_Y^k\\
\end{split}
\end{equation}




\subsection*{Diffusion of a non-conserved density: ($1,1,0$)-($0, 0,0$)}

\begin{equation}\label{diffusuz}
\begin{split}
TE^0 ={}& \dfrac{1}{2} A \, (\delta \mu)^2 + \dfrac{1}{2}\mathcal{K}\, \delta j_{k} \delta  j^k \\
TE^j ={}&  \delta \mu \, \delta j^j \\
T\sigma ={}& \Xi \, (\delta \mu)^2 + \mathcal{R} \, \delta j_{k} \delta  j^k\\
\end{split}
\end{equation}

\subsection*{Cattaneo model of diffusion: ($1,1,0$)-($1, 0,0$)}

\begin{equation}
\begin{split}
TE^0 ={}& \dfrac{1}{2} A \, (\delta \mu)^2 + \dfrac{1}{2}\mathcal{K}\, \delta j_{k} \delta  j^k \\
TE^j ={}&  \delta \mu \, \delta j^j \\
T\sigma ={}&  \mathcal{R} \, \delta j_{k} \delta  j^k\\
\end{split}
\end{equation}

\subsection*{Barotropic perfect fluid: ($1,1,0$)-($1, 1,0$)}

\begin{equation}
\begin{split}
TE^0 ={}& \dfrac{1}{2} \dfrac{n^2 (\delta \mu)^2 }{(\rho {+} P)c_s^2}+ \dfrac{1}{2} (\rho {+} P) \, \delta u_{k} \delta  u^k \\
TE^j ={}&  n\delta \mu \, \delta u^j \\
T\sigma ={}& 0\\
\end{split}
\end{equation}


\subsection*{Bulk viscous fluid at $\mu=0$: ($2,1,0$)-($1, 1,0$)}

\begin{equation}
\begin{split}
TE^0 ={}& \dfrac{1}{2} \bigg[ \dfrac{c_v}{T} (\delta T)^2 + b_0 (\delta \Pi)^2 + (\rho {+} P)\, \delta u_{k} \delta  u^k \bigg] \\
TE^j ={}&  (s \delta T +\delta \Pi) \delta u^j \\
T\sigma ={}& \dfrac{(\delta \Pi)^2}{\zeta} \\
\end{split}
\end{equation}


\subsection*{Perfect fluid: ($2,1,0$)-($2, 1,0$)}

\begin{equation}
\begin{split}
TE^0 ={}& \dfrac{1}{2} \bigg[ P_{TT} (\delta T)^2 +2 P_{T\mu} \delta T \delta \mu + P_{\mu \mu} (\delta \mu)^2 + (\rho {+} P)\, \delta u_{k} \delta  u^k \bigg] \\
TE^j ={}&  (s \delta T +n \delta \mu) \delta u^j \\
T\sigma ={}& 0 \\
\end{split}
\end{equation}


\subsection*{Coupled diffusion of two conserved densities: ($2,2,0$)-($2, 0,0$)}

\begin{equation}
\begin{split}
TE^0 ={}& \dfrac{1}{2} \bigg[ P_{TT} (\delta T)^2 +2 P_{T\mu} \delta T \delta \mu + P_{\mu \mu} (\delta \mu)^2 + \mathcal{K}^{ss} \delta s^k \delta s_{k} + 2\mathcal{K}^{sn} \delta s^k \delta n_{k} + \mathcal{K}^{nn} \delta n^k \delta n_{k} \bigg] \\
TE^j ={}&  \delta T \, \delta s^j + \delta \mu \, \delta n^j  \\
T\sigma ={}& \mathcal{R}_{ss} \delta s^k \delta s_{k} + 2\mathcal{R}_{sn} \delta s^k \delta n_{k} + \mathcal{R}_{nn} \delta n^k \delta n_{k} \\
\end{split}
\end{equation}


\subsection*{Heat conductive bulk viscous fluid at $\mu=0$: ($2,2,0$)-($1, 1,0$)}

\begin{equation}
\begin{split}
TE^0 ={}& \dfrac{1}{2} \bigg[ \dfrac{c_v}{T} (\delta T)^2 + b_0 (\delta \Pi)^2 + (\rho {+} P)\, \delta u_{k} \delta  u^k + 2\delta u_k \delta q^k + b_1 \delta q_k \delta q^k \bigg]  \\
TE^j ={}&  (s \delta T +\delta \Pi) \delta u^j + (T^{-1}\delta T -a_0 \delta \Pi) \delta q^j \\
T\sigma ={}& \dfrac{(\delta \Pi)^2}{\zeta} + \dfrac{\delta q_k \delta q^k}{\kappa T} \\
\end{split}
\end{equation}

\subsection*{Heat conductive fluid at $\mu \neq 0$: ($2,2,0$)-($2, 1,0$)}

\begin{equation}
\begin{split}
TE^0 ={}& \dfrac{1}{2} \bigg[ P_{TT} (\delta T)^2 +2 P_{T\mu} \delta T \delta \mu + P_{\mu \mu} (\delta \mu)^2 + (\rho {+} P)\, \delta u_{k} \delta  u^k + 2\delta u_k \delta q^k + b_1 \delta q_k \delta q^k \bigg]  \\
TE^j ={}&  (s \delta T +n \delta \mu) \delta u^j + T^{-1}\delta T \delta q^j \\
T\sigma ={}&  \dfrac{\delta q^k \delta q_k}{\kappa T} \\
\end{split}
\end{equation}

\subsection*{Relativistic superfluid: ($2,2,0$)-($2, 2,0$)}

\begin{equation}
\begin{split}
TE^0 ={}& \dfrac{1}{2} \bigg[ P_{TT} (\delta T)^2 +2 P_{T\mu} \delta T \delta \mu + P_{\mu \mu} (\delta \mu)^2 + \mathcal{K}^{ss} \delta s^k \delta s_{k} + 2\mathcal{K}^{sn} \delta s^k \delta n_{k} + \mathcal{K}^{nn} \delta n^k \delta n_{k} \bigg] \\
TE^j ={}&  \delta T \, \delta s^j + \delta \mu \, \delta n^j  \\
T\sigma ={}&0\\
\end{split}
\end{equation}

\subsection*{Maxwell material at $\mu=0$: ($1,1,1$)-($1, 1,0$)}

\begin{equation}
\begin{split}
TE^0 ={}& \dfrac{1}{2} \bigg[  \dfrac{c_v}{T} (\delta T)^2+ (\rho {+} P)\, \delta u_{k} \delta  u^k + b_2 \delta \Pi^{jk} \delta \Pi_{jk} \bigg] \\
TE^j ={}& s \delta T \delta u^j + \delta \Pi^{jk} \delta u_k \\
T\sigma ={}& \dfrac{\delta \Pi^{jk} \delta \Pi_{jk}}{2\eta} \\
\end{split}
\end{equation}


\subsection*{Elastic material at $\mu=0$: ($1,1,1$)-($1, 1,1$)}

\begin{equation}
\begin{split}
TE^0 ={}& \dfrac{1}{2} \bigg[  \dfrac{c_v}{T} (\delta T)^2+ (\rho {+} P)\, \delta u_{k} \delta  u^k + \dfrac{ \delta \Pi^{jk} \delta \Pi_{jk}}{2G} \bigg] \\
TE^j ={}& s \delta T \delta u^j + \delta \Pi^{jk} \delta u_k \\
T\sigma ={}& 0 \\
\end{split}
\end{equation}


\subsection*{Burgers material at $\mu=0$: ($1,1,2$)-($1, 1,0$)}

\begin{equation}
\begin{split}
TE^0 ={}& \dfrac{1}{2} \bigg[  \dfrac{c_v}{T} (\delta T)^2+ (\rho {+} P)\, \delta u_{k} \delta  u^k + b_2 \delta \Pi^{jk} \delta \Pi_{jk}+b_2 \delta \Lambda^{jk} \delta \Lambda_{jk}\bigg] \\
TE^j ={}& s \delta T \delta u^j + \delta \Pi^{jk} \delta u_k \\
T\sigma ={}& \xi_1 \delta \Pi^{kl}\delta \Pi_{kl}+2\xi_2 \delta \Pi^{kl}\delta \Lambda_{kl}+\xi_3\delta \Lambda^{kl}\delta \Lambda_{kl} \\
\end{split}
\end{equation}


\subsection*{Israel-Stewart theory in a general frame at $\mu =0$: ($3,2,1$)-($1, 1,0$)}

\begin{equation}
\begin{split}
TE^0 ={}& \dfrac{1}{2} \bigg[  \dfrac{c_v}{T} (\delta T)^2+ 2\dfrac{\delta \mathcal{A} \delta T}{T}+b_{1} (\delta \mathcal{A})^2+2b_2 \delta \Pi \delta \mathcal{A}+ b_3 (\delta \Pi)^2+(\rho {+} P)\, \delta u^{k} \delta  u_k + 2 \delta u^k \delta q_k+ b_4 \delta q^k \delta q_k + b_5 \delta \Pi^{jk} \delta \Pi_{jk} \bigg] \\
TE^j ={}& (s \delta T +\delta \Pi)\delta u^j  +(T^{-1} \delta T -a_{1} \delta \mathcal{A}- a_2 \delta \Pi )\delta q^j + (\delta u_k -a_3 \delta q_k)\delta \Pi^{jk}  \\
T\sigma ={}& \xi_1 (\delta \mathcal{A})^2 + 2\xi_2 \, \delta \mathcal{A} \delta \Pi +\xi_3 (\delta \Pi)^2 + \dfrac{\delta q^k \delta q_k}{\kappa T} +\dfrac{\delta \Pi^{jk} \delta \Pi_{jk}}{2\eta} \\
\end{split}
\end{equation}

\subsection*{Israel-Stewart theory at $\mu \neq 0$: ($3,2,1$)-($2, 1,0$)}

\begin{equation}
\begin{split}
TE^0={}& \dfrac{1}{2} \bigg[P_{TT} (\delta T)^2 +2 P_{T\mu} \delta T \delta \mu + P_{\mu \mu} (\delta \mu)^2  + (\rho{+}P) \delta u^k \delta u_k + 2 \delta u^k \delta q_k + b_1 \delta q^k \delta q_k +b_0 (\delta \Pi)^2 + b_2 \delta \Pi^{jk} \delta \Pi_{jk} \bigg] \\
TE^j={}& (s\delta T+ n\delta \mu +\delta \Pi)\delta u^j +(T^{-1} \delta T -a_0 \delta \Pi)\delta q^j + ( \delta u_k -a_1 \delta q_k)\delta \Pi^{jk}  \\
T\sigma ={}&  \dfrac{(\delta \Pi)^2}{\zeta} + \dfrac{\delta q^k \delta q_k}{\kappa T} +\dfrac{\delta \Pi^{jk} \delta \Pi_{jk}}{2\eta}  \\
\end{split}
\end{equation}


\subsection*{Elastic heat conducting material: ($3,2,1$)-($3, 1,1$)}
\begin{equation}
\begin{split}
TE^0={}& \dfrac{1}{2} \bigg[P_{TT} (\delta T)^2 +2 P_{T\mu} \delta T \delta \mu + P_{\mu \mu} (\delta \mu)^2   + (\rho{+}P) \delta u^k \delta u_k + 2 \delta u^k \delta q_k + b_1 \delta q^k \delta q_k + \dfrac{(\delta \Pi)^2}{K}+ \dfrac{\delta \Pi^{jk} \delta \Pi_{jk}}{2G} \bigg]   \\
TE^j={}& (s\delta T +n\delta \mu+ \delta \Pi)\delta u^j +(T^{-1} \delta T -a_0 \delta \Pi)\delta q^j + ( \delta u_k -a_1 \delta q_k)\delta \Pi^{jk}   \\
T\sigma ={}&  \dfrac{\delta q^k \delta q_k}{\kappa T}  \\
\end{split}
\end{equation}

\subsection*{Isotropic supersolid\footnote{Note that in the Andreev-Lifshitz theory for supersolid hydrodynamics \cite{Andreev1969,Sears2010} the coefficients $\mathcal{Z}_1$ and $\mathcal{Z}_2$ are set to zero. Here, to keep the class ``universal'', we have included them, because they are allowed by symmetry.}: ($3,2,1$)-($3, 2,1$)}


\begin{equation}\label{infosupsol}
\begin{split}
TE^0 ={}& \dfrac{1}{2} \bigg[P_{TT} (\delta T)^2 +2 P_{T\mu} \delta T \delta \mu + P_{\mu \mu} (\delta \mu)^2   + \mathcal{K}^{ss}\delta s^k \delta s_k + 2\mathcal{K}^{sn}\delta s^k \delta n_k + \mathcal{K}^{nn}\delta n^k \delta n_k + \dfrac{(\delta \Pi)^2}{K} {+} \dfrac{\delta \Pi^{kl}\delta \Pi_{kl}}{2G} \bigg]   \\
TE^j ={}& \delta T \delta s^j +\delta \mu \delta n^j + (s^{-1} \delta s^j-\mathcal{Z}_1 \delta n^j) \delta \Pi  + (s^{-1} \delta s_k -\mathcal{Z}_2 \delta n_k) \delta \Pi^{jk}    \\
T\sigma ={}& 0  \\
\end{split}
\end{equation}







\label{lastpage}

\end{document}
% ``






%The integral of each conserved current $j_I^\mu$ over a Cauchy surface $\Sigma$ defines the corresponding conserved charge,
% \begin{equation}
%     Q_I = \int_\Sigma j_I^\mu d\Sigma_\mu \, ,
% \end{equation}
% whose value is constant, and does not depend on the Cauchy surface. Similarly to what we did with the fields $\varphi^A$, also the charges $Q_I$ can be decomposed into irreducible tensors, and classified depending on how they transform under rotations. For example, energy and particle number are rotational scalars, while momentum and superfluid winding number \cite{Termo} are rotational vectors. This produces a second string of integers $(\bar{\mathfrak{g}}_1,\bar{\mathfrak{g}}_2,\bar{\mathfrak{g}}_3,...)$, where $\bar{\mathfrak{g}}_n$ counts how many charges belong to the $n-$th irreducible tensor category. This is the second ingredient of our classification.

In the linear regime, an isotropic supersolid can be interpreted as a system of class $(3,2,1)-(3,2,1)$, whose fields are $\mu^I=\{\delta T, \delta \mu, \delta \Pi, \delta s^k, \delta n^k, \delta \Pi^{kl} \}$ which may be interpreted as the perturbations to respectively temperature, chemical potential, elastic bulk stress, entropy flux, particle flux, and elastic shear stress. No dissipative effects are considered here (i.e. there are as many fields as conservation laws). Thus, the supersolid information current can be easily guessed by combining the superfluid information current \eqref{infosup} with the elastic information current \eqref{infoelas}, giving
% \begin{equation}\label{infosupsol}
% \begin{split}
% TE^0 ={}& \dfrac{1}{2} \bigg[\dfrac{\partial s}{\partial T}\bigg|_{\mu}  \! \! \! (\delta T)^2 {+} 2 \dfrac{\partial s}{\partial \mu}\bigg|_{T}  \! \! \! \delta T \, \delta \mu {+} \dfrac{\partial n}{\partial \mu}\bigg|_{T}  \! \! \!(\delta \mu)^2 + \mathcal{K}^{ss}\delta s^k \delta s_k + 2\mathcal{K}^{sn}\delta s^k \delta n_k + \mathcal{K}^{nn}\delta n^k \delta n_k + \dfrac{(\delta \Pi)^2}{K} {+} \dfrac{\delta \Pi^{kl}\delta \Pi_{kl}}{2G} \bigg]  , \\
% TE^j ={}& \delta T \delta s^j +\delta \mu \delta n^j + s^{-1} \delta \Pi \delta s^j + s^{-1} \delta \Pi^{jk} \delta s_k \, , \\
% T\sigma ={}& 0 \, , \\
% \end{split}
% \end{equation}
% where $s$ is the entropy density, $n$ is the particle density, $\mathcal{K}^{XY}$ is the entrainment matrix (see section \ref{entrainment}), $K$ is the bulk modulus, and $G$ is the shear modulus. In the formula for $TE^j$, we postulated the coupling $ s^{-1} \delta \Pi \delta s^j + s^{-1} \delta \Pi^{jk} \delta s_k$, because only the normal component participates in the elastic behavior, and $\delta u_N^j:=\delta s^j/s$ is (by definition) the velocity of the normal flow. Note that, in the fully non-linear supersolid theory, the elastic stress $\Pi^{kl} +\Pi \, \delta^{kl}$ is not symmetric. However, it becomes symmetric in the linear regime \cite{Andreev1969}, so that (as usual) $\delta \Pi^{kl}$ is assumed symmetric and traceless. 
% The field equations \eqref{finally}, computed from \eqref{infosupsol}, read 
% \begin{equation}\label{supersoliffieldeq}
% \begin{split}
% \partial_t \bigg[\dfrac{\partial s}{\partial T}\bigg|_{\mu}   \delta T + \dfrac{\partial s}{\partial \mu}\bigg|_{T}   \delta \mu  \bigg] +\partial_j \delta s^j ={}& 0 \, , \\
% \partial_t \bigg[\dfrac{\partial n}{\partial \mu}\bigg|_{T}  \delta \mu + \dfrac{\partial s}{\partial \mu}\bigg|_{T}   \delta T  \bigg] +\partial_j \delta n^j={}& 0 \, , \\
% \partial_t (\mathcal{K}^{ss} \delta s_k + \mathcal{K}^{sn}\delta n_k) +\partial_k (\delta T+s^{-1} \delta \Pi) + s^{-1} \partial_k \delta \Pi^j_k={}& 0 \, , \\
% \partial_t (\mathcal{K}^{nn} \delta n_k + \mathcal{K}^{sn}\delta s_k) +\partial_k \delta \mu={}& 0 \, , \\
% K^{-1} \partial_t \delta \Pi +\partial_j (s^{-1} \delta s^j)={}& 0 \, , \\
% (2G)^{-1} \partial_t \delta \Pi_{kl} + \braket{\partial_k (s^{-1}\delta s_l)}={}& 0 \, , \\
% \end{split}
% \end{equation}
% and they constitute a manifestly symmetric hyperbolic system. Our goal, now, is to show that this system is indeed the correct relativistic generalization of supersolid hydrodynamics \cite{Sears2010} (in the non-dissipative limit). 

% First, we see that, since by definition $\delta s^j=s\delta u_N^j$, the first equation of \eqref{supersoliffieldeq} simplifies to $\partial_t \delta s+\partial_j (s\delta u_N^j)=0$, i.e. the entropy is advected by the normal component, compare with eq.s (8) and (13) of \cite{Sears2010} (neglecting dissipation). Also, if we introduce the lattice displacement $\xi^k$, whose defining equation is $\partial_t \xi^k = \delta u_N^k$ (see eq.s (9) and (14) of \cite{Sears2010}), the last two equations of \eqref{supersoliffieldeq} can be integrated to give $\delta \Pi=-K\partial_j \xi^j$, and $\delta \Pi_{kl}=-2G\braket{\partial_k \xi_l}$, see eq.s (22) and (43) of \cite{Sears2010}. Furthermore, we know from Carter's multifluid theory \cite{cool1995} that the Landau superfluid velocity $\delta u_S^j$ is defined through the equation $\mu  \delta u_S^j = \mathcal{K}^{nn}\delta n^j+\mathcal{K}^{sn}\delta s^j$ (note that $\mu$ is the \textit{relativistic} chemical potential: It contains the ``$mc^2$'' term). As a consequence, the fourth equation of \eqref{supersoliffieldeq} becomes $\mu \partial_t \delta u_{Sk}+\partial_k \delta \mu=0$, which is the relativistic analogue of eq.s (12) and (16) of \cite{Sears2010}. One should keep in mind that, in the Newtonian limit, $\mu \rightarrow m$, where $m$ is the mass of the constituents. Additionally, note that what \citet{Sears2010} call ``$\mu$'' is actually $\mu/m -1$ for us.   


% Thanks to the well-known Maxwell relation
% \begin{equation}
%     \dfrac{\partial s}{\partial \mu}\bigg|_{T}  = \dfrac{\partial n}{\partial T}\bigg|_{\mu} \, ,
% \end{equation}
% the second equation of \eqref{supersoliffieldeq} reduces to the usual continuity equation, $\partial_t \delta n+\partial_j \delta n^j=0$, see eq. (10) of \cite{Sears2010}. However, one needs to be careful here, since in relativity the momentum density, $\delta T^{0j}$, and the current of rest mass, $m \delta n^j$, do not coincide. Indeed,  $\delta T^{0j}$ is must be a linear combination of both $\delta n^j$ and $\delta s^j \,$:
% \begin{equation}\label{totalmomentum}
%     \delta T^{0j} = (n\mathcal{K}^{nn}+s\mathcal{K}^{sn})\delta n^j +(n\mathcal{K}^{sn}+s\mathcal{K}^{ss})\delta s^j = \mu \delta n^j +T \delta s^j = \rho_S \delta u_S^j+\rho_N \delta u_N^j \, ,
% \end{equation}
% where $\rho_S$ is Landau's superfluid mass density, defined through equation \eqref{entratevoi!}, and $\rho_N$ is Landau's normal mass density, defined through the condition $sT+\mu n=\rho_S+\rho_N$ (relativistic partition of inertia \cite{GavassinoStabilityCarter2022}). The first equality sign in \eqref{totalmomentum} follows from multifluid hydrodynamics \cite{cool1995}, the second follows from \eqref{entratevoi!}, and the third follows from the definitions of $\delta u_S^j$ and $\delta u_N^j$. Combining the third and fourth equations of \eqref{supersoliffieldeq}, using the identity $\delta P =s\delta T +n \delta \mu$, we obtain $\partial_t \delta T^0_j+\partial_k (\delta P+\delta \Pi)+\partial_j\delta \Pi^j_k=0$, which matches eq.s (4), (11), and (15) of \cite{Sears2010}, thus completing our comparison.


% For example, take the case $\varphi^A=\{\delta n^0 , \delta n^j \}$. As we anticipated, the most general rotational invariant theory with these degrees of freedom must have the form
% \begin{equation}
% \begin{split}\label{gmebop}
%     E^0 ={}& \dfrac{1}{2} M_1 \, (\delta n^0)^2 + \dfrac{1}{2} M_2 \, \delta n^k \delta n_k \, , \\
%     E^j ={}& M_3 \, \delta n^0 \delta n^j \, , \\
%     \sigma ={}&  M_4 \, (\delta n^0)^2 +  M_5 \, \delta n^k \delta n_k \, , \\
%     \end{split}
% \end{equation}
% since no are terms are geometrically allowed. Here, $M_i$ are some free parameters, and they may be interpreted as the transport coefficients of the theory. Now, let us set $M_4=M_5=0$ (non-dissipative limit), and 
% \begin{equation}
%     M_1=M_2 \, c_s^2=M_3 >0 \, ,
% \end{equation}
% where $c_s^2>0$ is the speed of sound of the material under consideration. It is straightforward to verify the resulting formulas for $E^\mu$ and $\sigma$ are those of a barotropic perfect fluid \cite{GavassinoCausality2021}. Furthermore, if we compute \eqref{cinque} from \eqref{gmebop}, with the above choices of $M_i$, the resulting dynamical equations are indeed those of a barotropic perfect fluid. Thus, we can say that the barotropic perfect fluid is a ``particular realization'' of the general theory \eqref{gmebop}. 
% But there are other possible realizations. For example, take again the barotropic perfect fluid, and turn on dissipation through $M_5$. Then, if we define the charge diffusivity as $\kappa=M_2 c_s^2/M_5$, the system \eqref{cinque}, computed from \eqref{gmebop}, and divided for convenience by $M_2 c_s^2$, becomes
% \begin{equation}\label{ahkundt}
% \begin{split}
%     & \partial_\mu \delta n^\mu =0 \, , \\
%     & c_s^{-2} \partial_t \delta n_k +\partial_k \delta n^0 = -\kappa^{-1} \delta n_k \, . \\
%     \end{split}
% \end{equation}
% These are the equations of Cattaneo's model for diffusion \cite{cattaneo1958,Jou_Extended,rezzolla_book}, with relaxation time $\tau=\kappa/c_s^2$, and characteristic speed $c_s$. In the limit of small gradients, the second equation of \eqref{ahkundt} relaxes to Fick's law of diffusion, $\delta n_k=-\kappa \, \partial_k \delta n^0$, justifying the interpretation of $\kappa$ as the charge diffusivity.

% The above example reveals something very interesting






%In this way, the system of equations that arises from the triplet $\{\varphi^A,E^\mu,\sigma \}$, constructed as above, constitutes the most general (rotational invariant) theory that is consistent with the hypotheses of Theorems \ref{theo} and \ref{theo2}, and whose degrees of freedom $\varphi^A$ can be classified according to their geometric character consistently with the list $(\mathfrak{g}_1,\mathfrak{g}_2,\mathfrak{g}_3,...)$. Any theory with the same values of all $\mathfrak{g}_n$ must be a particular realization of this ``general theory'', $\{\varphi^A,E^\mu,\sigma \}$, for some specific choice of transport coefficients.

% This theorem establishes our notion of ``universality class''. In fact, it guarantees that, once we specify the transformation properties of the fields $\varphi^A$ under rotations (i.e. once we fix $\mathfrak{s},\mathfrak{v},...$), we can write the most general formulas for $E^\mu(\varphi^A)$ and $\sigma(\varphi^A)$ compatibly with rotational invariance (since the equilibrium state is assumed isotropic), and this determines the most general theory for such degrees of freedom. All other theories with the same geometric degrees of freedom must be  particular realizations of this general theory. For example, in the Supplemental material, we prove the following 
% \begin{theorem}\label{theo3}
% The linearised Israel-Stewart theory in the Eckart frame at finite chemical potential \cite{Hishcock1983} is the most general linear theory that satisfies the hypotheses of Theorem \ref{theo}, and with the following properties:
% \begin{itemize}
% \item The universality class is $(\mathfrak{s},\mathfrak{v},\mathfrak{t})=(3,2,1)$;
% \item Particles, energy, and momentum are conserved;
% \item The stress-energy tensor is symmetric.
% \end{itemize}
% \end{theorem}
% The above theorem tells us that, if the degrees of freedom of a hydrodynamic theory can be rearranged into something that ``looks like'' the degrees of the Israel-Stewart theory, then in the linear regime such theory becomes indistinguishable from the Israel-Stewart theory in the Eckart frame. This implies, for example, that in the Israel-Stewart framework all hydrodynamic frames are mathematically equivalent in the linear regime. Furthermore, we can also conclude that divergence-type theories \cite{Geroch1995}, Carter's theory for viscosity \cite{carter1991}, and the GENERIC theory \cite{Stricker2019} all become indistinguishable from the Israel-Stewart theory, near equilibrium. This interesting phenomenon has already been noticed in some specific cases \cite{Liu1986,PriouCOMPAR1991,GavassinoNonHydro2022,
% GavassinoGENERIC2022,SalazarZannias2022}. Now it is established in full generality. 

% In the following, we will give some examples of apparently different physical systems whose hydrodynamic descriptions become mathematically equivalent in the linear regime (as a result of Theorem \ref{theo2}). We say that such systems belong to the same ``universality class''. 

% 


% Now, the entropy production rate $\sigma$ is a non-negative definite quadratic form in $\varphi^A$, see equation \eqref{second}. Therefore, by the Sylvester law of inertia, we can always introduce an invertible change of variables, $\varphi^A=\mathcal{N}^A_C\tilde{\varphi}^C$, such that $\tilde{\sigma}_{CD}$ is diagonal, with diagonal entries equal to $0$ and $+1$. The fields $\tilde{\varphi}^C$ can also be ordered in such a way that $\tilde{\sigma}_{CD}$ takes the form of a block matrix:
% \begin{equation}\label{iknowthatIshold}
%     \tilde{\sigma}_{CD}=
% \left[\begin{array}{ c | c }
%     0 & 0 \\
%     \hline
%     0 & \mathbb{I}
%   \end{array}\right] ,
% \end{equation}
% where $\mathbb{I}$ is the identity matrix. On the other hand, we know from \eqref{conserviamodibrutto} that the null space of $\sigma$ is generated by all vectors of the form $\mathcal{N}_I^A Z^I$, for arbitrary $Z^I$. It follows that the size of the upper diagonal block in \eqref{iknowthatIshold} is determined by the number of conservation laws of the theory. Indeed, for each field $\tilde{\varphi}^C$ that ``falls'' into that block, the corresponding field equation is a conservation law. Hence, we can do something similar to what we did in step 1: We can classify the conservation laws of the system, by simply classifying the fields $\tilde{\varphi}^C$ for which $\tilde{\sigma}_{CC}=0$.












% \textit{Universality classes -} We are now ready to formulate a rigorous notion of ``universality class''. The idea is the following. Consider an arbitrary collection of fields $\varphi^A$. Classify such fields depending on how they transform under rotations (in the global equilibrium rest frame). For example, if among the fields there is a fourcurrent $\delta n^\mu$, then $\delta n^0$ should be counted as a rotational scalar, and $\delta n^j$ ($j=1,2,3$) should be counted as a rotational vector. Then, construct the most general expressions for $E^\mu$ and $\sigma$, as given in \eqref{second}, which are consistent with the transformation properties of the fields under rotations. In particular, since the equilibrium state is isotropic, the formulas for $E^\mu$ and $\sigma$ should not single out a preferred direction in space, and the resulting expressions should be formally invariant under rotations in the equilibrium global rest frame. For example, if our only degree of freedom is a fourcurrent $\delta n^\mu$, then $E^0$ and $\sigma$ can only contain a term proportional to $(\delta n^0)^2$ and a term proportional to $\delta n^k \delta n_k$, while $E^j$ can only contain a term proportional to $\delta n^0 \delta n^j$. In this way, the system of equations that arises from the triplet $\{\varphi^A,E^\mu,\sigma \}$, constructed as above, constitutes the most general (rotational invariant) theory that is consistent with the hypotheses of Theorems \ref{theo} and \ref{theo2}, and whose degrees of freedom are the fields $\varphi^A$. 

% But we can say even more. Let us group the fields $\varphi^A$ as irreducible tensors transforming under irreducible representations of the group of rotations $SO(3)$, and let us count how many fields there are for each irreducible tensor category. This allows us to assign to our theory a string of integers $(\mathfrak{g}_1,\mathfrak{g}_2,\mathfrak{g}_3,...)$, where $\mathfrak{g}_n$ is the number of fields that belong to the $n-$th irreducible tensor category. In practice, we may fix $\mathfrak{g}_1=\mathfrak{s}$ to be the number of scalars, $\mathfrak{g}_2=\mathfrak{v}$ the number of vectors, $\mathfrak{g}_3=\mathfrak{t}$ the number of symmetric traceless tensors with two indices, and so on. It is evident that our previously constructed theory $\{ \varphi^A,E^\mu,\sigma \}$ is the most general rotational invariant theory that one can build 



% Now, consider another rotational invariant theory, $\{\tilde{\varphi}^C, \tilde{E}^\mu,\tilde{\sigma} \}$, with the same values of all $\mathfrak{g}_n$. Clearly, we can always invent a one-to-one correspondence, $\varphi^A=\mathcal{N}^A_C \tilde{\varphi}^C$, which sends scalars into scalars, vectors into vectors, and, more in general, tensors belonging to the $n-$th irreducible tensor category into tensors belonging to the same category. This mapping preserves rotational invariance, and it allows us to express $E^\mu$ and $\sigma$ is terms of $\tilde{\varphi}^C$. But, since $E^\mu(\varphi^A)$ and $\sigma(\varphi^A)$ are most general information current and entropy production rate that we can construct compatibly with rotational invariance, there must be a specific choice of coefficients for which equation \eqref{potato} is valid. Therefore, by application of Theorem \ref{theo2}, we can conclude that, for such choice of coefficients, $\{\tilde{\varphi}^C, \tilde{E}^\mu,\tilde{\sigma} \}$ is equivalent to $\{\varphi^A,E^\mu,\sigma \}$, which means that any theory with exactly $\mathfrak{g}_n$ irreducible tensors belong to the $n-$th category (for all $n$) must a particular case of the ``universal theory'' $\{\varphi^A,E^\mu,\sigma \}$, where we fix the parameters to have some specific value.


% To have a first taste of how this works, consider Cattaneo's model for heat conduction \cite{GavassinoNonHydro2022}, which belongs to the class $(1,1,0,0,0...)$:
% \begin{equation}\label{cvuno}
%     \begin{split}
%         &  c_v \partial_t \delta T +\partial_j \delta q^j =0 \, ,\\
%         & T\beta_1 \partial_t \delta q_j + \partial_j \delta T = -\kappa^{-1} \delta q_j \, ,
%     \end{split}
% \end{equation}
% where $T$, $q_j$, $c_v$, $\kappa$ and $\beta_1$ are temperature, heat flux, heat capacity per unit volume, heat conductivity, and a second order transport coefficient. If we set $\kappa^{-1}=0$, and $\beta_1=(Ts)^{-1}$, where $s$ is the entropy density, the system \eqref{cvuno} becomes
% \begin{equation}\label{cvdue}
%     \begin{split}
%         &  c_v \partial_t \delta T +Ts \, \partial_j \delta u^j =0 \, ,\\
%         & T \partial_t \delta u_j + \partial_j \delta T =0 \, ,
%     \end{split}
% \end{equation}
% with $\delta q^j=Ts \, \delta u^j$. The above equations describe a perfect fluid at zero chemical potential, which also belongs to the class $(1,1,0,0,0...)$. As we can see, the perfect fluid is indistinguishable from Cattaneo's heat conduction (in the linear regime), with $\kappa=+\infty$, and this follows from the fact that the degrees of freedom have the same geometric character. 

%
%Let us quickly review some basics of relativistic thermodynamics. Consider a fluid in contact with a heat bath ``$H$'', and assume that the total system ``fluid$+$bath'' is isolated. Then, in an arbitrary process
%\begin{equation}
%\Delta S+\Delta S_H \geq 0 \, , \spc \Delta Q^I_H =-\Delta Q^I  \, .
%\end{equation}
%



%The equilibrium state maximizes the entropy $S$ for fixed values of all conserved charges $Q^I$. Then, chosen an equilibrium state, there must be some constants $\alpha^\star_I$ (which identify that specific state) such that
%\begin{equation}
%dS +\alpha^\star_I dQ^I=0 \, .
%\end{equation}


%The state of thermodynamic equilibrium is always the maximum or minimum of some thermodynamic potential $\Phi$ (e.g., the entropy or the free energy). Close to equilibrium, such potential can be expressed as $\Phi=\Phi_\text{eq}+\Delta \Phi$, with
%\begin{equation}
%\Phi(\Sigma) = c \int_\Sigma E^\mu \, d\Sigma_\mu \, ,
%\end{equation}
%where $c$ is a constant that depends on the potential we are choosing, and $\Sigma$ is Cauchy surface. The vector field $E^\mu$ is known as ``information current'', and it is timelike
%
%Suppose that the linearised field equations are of first order:
%\begin{equation}\label{prima}
%M^\mu_{AB}\partial_\mu \varphi^B = -\Xi_{AB} \varphi^B \, ,
%\end{equation}
%where $M^0_{AB}$ is an invertible matrix. If we contract both sides with $\varphi^A$, and bring all terms on the left-hand side, we obtain
%\begin{equation}\label{prima2}
%\varphi^A M^\mu_{AB}\partial_\mu \varphi^B +\varphi^A \Xi_{AB} \varphi^B=0 \, .
%\end{equation}
%Suppose also that there is a quadratic information current, whose four-divergence is a quadratic form in the fields:
%\begin{equation}\label{second}
%\begin{split}
%& \partial_\mu E^\mu = -\sigma  \, , \\
%& E^\mu = \dfrac{1}{2} E^\mu_{AB} \varphi^A \varphi^B \, ,  \\
%& \sigma = \dfrac{1}{2} \sigma_{AB} \varphi^A \varphi^B \, . \\
%\end{split}
%\end{equation}
%The matrices $E^\mu_{AB}$ and $\sigma_{AB}$ are symmetric. Stability and causality imply that $E^\mu$ is time-like future-directed. Hence, also $E^0_{AB}$ is invertible. Combining the equation in \eqref{second}, we obtain (bringing all terms on the left)
%\begin{equation}\label{terza}
%\varphi^A E^\mu_{AB} \partial_\mu \varphi^B + \varphi^A \dfrac{\sigma_{AB}}{2} \varphi^B =0 \, .
%\end{equation}
%We always have the freedom to multiply both sides of \eqref{prima} by an invertible matrix $\mathcal{N}^A_C$, and this does not change the dynamics of the system. Hence, I can always ``fix'' the matrix $M^0_{AB}$ to coincide with $E^0_{AB}$ (since they are both invertible). But then, if we subtract \eqref{prima2} to \eqref{terza}, the terms $\varphi^A E^0_{AB} \partial_t \varphi^B$ cancel out, and we are left with
%\begin{equation}
%\varphi^A(E^j_{AB}-M^j_{AB})\partial_j \varphi^B +\varphi^A \bigg( \dfrac{\sigma_{AB}}{2} - \Xi_{AB} \bigg) \varphi^B =0 \, .
%\end{equation}
%But this must be true at any point, and for any possible choice of initial data. In particular, if we choose some initial data with no spatial gradients, the first term vanishes identically, and we obtain the condition
%\begin{equation}
%\Xi_{AB}= \dfrac{1}{2} \sigma_{AB} + \Xi_{[AB]} \, .
%\end{equation}
%But this implies that $\varphi^A(E^j_{AB}-M^j_{AB})\partial_j \varphi^B$ must be zero for any choice of $\varphi^A$ and for any choice of $\partial_j \varphi^B$, hence
%\begin{equation}
%M^j_{AB} = E^j_{AB} \, .
%\end{equation}
%Recalling that we have also set $M^0_{AB}=E^0_{AB}$, we can equivalently rewrite equation \eqref{prima} in the form
%\begin{equation}
%\partial_\mu \dfrac{\partial E^\mu}{\partial \varphi^A} = -\dfrac{1}{2} \dfrac{\partial \sigma}{\partial \varphi^A} -\Xi_{[AB]} \varphi^B \, .
%\end{equation}
%This proves that the linearised Israel-Stewart theory is symmetric-hyperbolic (provided that it is stable) in any possible hydrodynamic frame, also the ``general frame'', and with an arbitrary number of currents. Furthermore, its equations can always be computed from the information current as above, which is just the method I had formulated in my previous papers!







%\textit{Chemistry is equivalent to bulk viscosity -} Consider a fluid mixture of two species, where one (say, baryons, $b$) has a conserved current, and the other (say, neutral pions, $\uppi^0$) does not have a conserved current. Mixtures of this kind are frequent in neutron stars \cite{Migdal1973,Suzuki1973,Haensel1982}. The first law of thermodynamics reads $d\rho= Tds+\mu db-\mathbb{A}d\uppi$, where $\rho$, $s$, $b$ and $\uppi$ are densities of energy, entropy, baryons, and pions, $T$ is the temperature, $\mu$ and $-\mathbb{A}$ are respectively the baryon and the pion chemical potential. $\mathbb{A}$ is also the affinity \cite{PrigoginebookModernThermodynamics2014} of reactions like
% \begin{equation}
% b+b \ce{ <=> } b+b+\uppi^0 \, ,
% \end{equation}
% and it must vanish in thermodynamic equilibrium. Working in the equilibrium rest frame, the information current and entropy production rate for this system are \cite{GavassinoGibbs2021}
% \begin{equation}\label{infochemical}
% \begin{split}
% TE^0 ={}& \dfrac{1}{2} \bigg[\dfrac{\partial T}{\partial s}\bigg|_{b,\mathbb{A}} \! \! \! (\delta s)^2 + 2 \dfrac{\partial T}{\partial b}\bigg|_{s,\mathbb{A}} \! \! \! \delta s \, \delta b + \dfrac{\partial \mu}{\partial b}\bigg|_{s,\mathbb{A}} \! \! \!(\delta b)^2 \\
% & -\dfrac{\partial \uppi}{\partial \mathbb{A}}\bigg|_{s,b} (\delta \mathbb{A})^2 + (\rho {+} P) \delta u^k \delta u_k  \bigg] \, , \\
% TE^j ={}& \bigg[  \dfrac{\partial P}{\partial s}\bigg|_{b,\mathbb{A}} \delta s+ \dfrac{\partial P}{\partial b}\bigg|_{s,\mathbb{A}} \delta b+ \dfrac{\partial P}{\partial \mathbb{A}}\bigg|_{s,b} \delta \mathbb{A} \bigg] \delta u^k \, , \\
% T\sigma ={}& \Xi \, (\delta \mathbb{A})^2 \, , \\
% \end{split}
% \end{equation}
% where $P(s,b,\mathbb{A})$ is the total non-equilibrium pressure and $\Xi$ is the net pion production rate. All background quantities are evaluated at $\mathbb{A}=0$, which is the condition for chemical equilibrium. One can easily verify that, if we set $\varphi^A=\{\delta s, \delta b, \delta \mathbb{A}, \delta u^k \}$, and we write the system \eqref{cinque} explicitly, we obtain a system of field equations that is indeed equivalent to the linearised equations of motion of the mixture \cite{BulkGavassino}, but rewritten in a manifestly symmetric-hyperbolic form. Let us now introduce the following notation:
% \begin{equation}
% \begin{split}
% \delta \Pi ={}& \dfrac{\partial P}{\partial \mathbb{A}}\bigg|_{s,b} \delta \mathbb{A} \, , \\
% \zeta ={}& \dfrac{1}{\Xi} \bigg( \dfrac{\partial P}{\partial \mathbb{A}}\bigg|_{s,b} \bigg)^2 \, , \\
% \tau_\Pi ={}& -\dfrac{1}{\Xi} \dfrac{\partial \uppi}{\partial \mathbb{A}} \bigg|_{s,b} \, . \\
% \end{split}
% \end{equation}
% Then, equation \eqref{infochemical} becomes
% \begin{equation}\label{infobulk}
% \begin{split}
% TE^0 ={}& \dfrac{1}{2} \bigg[\dfrac{\partial T}{\partial s}\bigg|_{b}  \! (\delta s)^2 + 2 \dfrac{\partial T}{\partial b}\bigg|_{s}  \! \delta s \, \delta b + \dfrac{\partial \mu}{\partial b}\bigg|_{s} \!(\delta b)^2 \\
% & \dfrac{\tau_\Pi}{\zeta} (\delta \Pi)^2 + (\rho {+} P) \delta u^k \delta u_k  \bigg] \, , \\
% TE^j ={}& \bigg[  \dfrac{\partial P}{\partial s}\bigg|_{b} \delta s+ \dfrac{\partial P}{\partial b}\bigg|_{s} \delta b \bigg] \delta u^k + \delta \Pi \, \delta u^k \, , \\
% T\sigma ={}& (\delta \Pi)^2/\zeta \, . \\
% \end{split}
% \end{equation}
% This are the information current and entropy production rate of a bulk-viscous baryonic fluid, modelled using the Israel-Stewart theory \cite{Hishcock1983}, with bulk visconsity $\zeta$, relaxation time $\tau_\Pi$, and bulk-viscous stress $\delta \Pi$. In moving from equation \eqref{infochemical} to \eqref{infobulk}, we changed notation,
% \begin{equation}
% \dfrac{\partial f}{\partial g}\bigg|_{h,\mathbb{A}} = \dfrac{\partial f}{\partial g}\bigg|_{h} \, ,  
% \end{equation} 
% because the Israel-Stewart equilibrium equation of state $\rho(s,b)$ is obtained by imposing the constraint $\mathbb{A}=0$ in the non-equilibrium equation of state $\rho(s,b,\uppi)$, meaning that the Israel-Stewart thermodynamic derivatives are inevitably partial derivatives at constant $\mathbb{A}$ (with $\mathbb{A}=0$). 

% We have explicitly constructed a mapping from the baryon-pion mixture to the Israel-Stewart theory for bulk viscosity, which satisfies \eqref{potato}. Hence, Theorem \ref{theo2} guarantees that, in the linear regime, the baryon-pion mixture and Israel-Stewart bulk viscosity are mathematically equivalent. Indeed, this equivalence has already been observed, both theoretically \cite{BulkGavassino,GavassinoRadiazione,Camelio2022} and numerically \cite{CamelioSimulations2022}, but it had never been established with this level of mathematical rigour. In fact, not only these two systems ``contain the same physics''. In the linear regime, they are really the same system of equations. This implies that if we know the conditions for stability and causality of Israel-Stewart bulk viscosity, we also know those of the mixture. For example, the characteristic speed of Israel-Stewart bulk viscosity is \cite{Causality_bulk}
% \begin{equation}
% \begin{split}
% c_{UV}^2 ={}& \dfrac{\partial P}{\partial \rho} \bigg|_{s/b} \!\! + \dfrac{\zeta}{\tau_\Pi (\rho {+} P)} = \\
% & \dfrac{\partial P}{\partial \rho} \bigg|_{s/b,\mathbb{A}} \!\! \!\! -  \dfrac{\partial P}{\partial \mathbb{A}}\bigg|_{s,b} \dfrac{\partial P}{\partial \uppi}\bigg|_{s,b} \, \dfrac{1}{\rho {+} P} = \\
% & \dfrac{\partial P}{\partial \rho} \bigg|_{s/b,\mathbb{A}} \!\! \!\! +  \dfrac{\partial P}{\partial \mathbb{A}}\bigg|_{\rho,s/b} \dfrac{\partial \mathbb{A}}{\partial \rho}\bigg|_{s/b,\uppi/b} \!\!= \dfrac{\partial P}{\partial \rho}\bigg|_{s/b,\uppi/b} \, , \\
% \end{split}
% \end{equation}
% which is, indeed, the characteristic speed of the mixture. In conclusion, we can say that the baryon-pion mixture belongs to the ``bulk viscosity universality class''.

% \textit{Superfluidity is equivalent to heat conduction -} The relativistic version of Landau's two-fluid model  for superfluidity \cite{landau6,khalatnikov_book} is due to Carter \cite{carter1991,cool1995,langlois98} and Khalatnikov \cite{lebedev1982,Carter_starting_point,carter92}, see also \cite{Son2001,Gusakov2007,andersson2007review,Termo,
% GavassinoIordanskii2021}. In the linear regime, the fields of the Carter-Khalatikov theory are $\varphi^A=\{ \delta T, \delta \mu, \delta s^k, \delta n^k \}$, namely temperature, chemical potential, entropy flux, and particle flux. The information current and entropy production rate are \cite{GavassinoStabilityCarter2022} 
% \begin{equation}\label{infosup}
% \begin{split}
% TE^0 ={}& \dfrac{1}{2} \bigg[\dfrac{\partial s}{\partial T}\bigg|_{\mu}  \! (\delta T)^2 + 2 \dfrac{\partial s}{\partial \mu}\bigg|_{T}  \! \delta T \, \delta \mu + \dfrac{\partial n}{\partial \mu}\bigg|_{T} \!(\delta \mu)^2 \\
% & \mathcal{K}^{ss}\delta s^k \delta s_k + 2\mathcal{K}^{sn}\delta s^k \delta n_k + \mathcal{K}^{nn}\delta n^k \delta n_k  \bigg] \, , \\
% TE^j ={}& \delta T \delta s^j +\delta \mu \delta n^j \, , \\
% T\sigma ={}& 0 \, , \\
% \end{split}
% \end{equation}
% where $s$ and $n$ are the entropy and particle density. The phenomenological coefficients $\mathcal{K}^{XY}$ constitute the so-called ``entrainment matrix'', and they are related to the Landau ``superfluid mass density'' $\rho_S$ \cite{landau9} as follows:
% \begin{equation}
% \mathcal{K}^{XY}= 
% \begin{bmatrix}
%   \dfrac{1}{s^2}\bigg(Ts-\mu n + \dfrac{\mu^2n^2}{\rho_S}\bigg) & \, \, \dfrac{\mu}{s} \bigg( 1-\dfrac{\mu n}{\rho_S} \bigg)  \\
%   \dfrac{\mu}{s} \bigg( 1-\dfrac{\mu n}{\rho_S} \bigg) &  \dfrac{\mu^2}{\rho_S} \\
% \end{bmatrix} .
% \end{equation}
% Now, let us introduce the following change of notation:
% \begin{equation}\label{beta1}
% \begin{split}
% \delta u^k ={}& \delta n^k/n \, ,  \\
% \delta q^k ={}& T(\delta s^k-s\delta u^k) \, , \\
% \beta_1 ={}&  \dfrac{1}{T^2 s^2}\bigg(Ts-\mu n + \dfrac{\mu^2n^2}{\rho_S}\bigg) \, , \\
% \end{split}
% \end{equation}
% where $\delta u^k$ can be interpreted as the flow velocity (in the Eckart frame \cite{Eckart40}), and $\delta q^k$ can be interpreted as the heat flux.
% Then, equation \eqref{infosup} becomes
% \begin{equation}\label{infoheat}
% \begin{split}
% TE^0 ={}& \dfrac{1}{2} \bigg[\dfrac{\partial s}{\partial T}\bigg|_{\mu}  \! (\delta T)^2 + 2 \dfrac{\partial s}{\partial \mu}\bigg|_{T}  \! \delta T \, \delta \mu + \dfrac{\partial n}{\partial \mu}\bigg|_{T} \!(\delta \mu)^2 \\
% & (\rho {+} P)\delta u^k \delta u_k +2\delta u^k \delta q_k + \beta_1 \delta q^k \delta q_k  \bigg] \, , \\
% TE^j ={}& (s  \delta T +n  \delta \mu)\delta u^j + \delta T \delta q^j/T \, , \\
% T\sigma ={}& \delta q^k \delta q_k /\kappa T^2 \spc  (\text{with }\kappa=+\infty) \, . \\
% \end{split}
% \end{equation}
% We have recovered the information current and entropy production rate of the Israel-Stewart theory for heat conduction \cite{OlsonRegularCarter1990}, with infinite heat conductivity $\kappa$. Then, we can invoke Theorem \ref{theo2}, and we arrive at a surprisingly simple and intuitive result: in the linear regime, a relativistic superfluid is described by exactly the same equations of an ordinary heat-conducting fluid, in the limit of infinite heat conductivity. In other words, Landau's two-fluid model belongs to the ``heat conduction universality class''. Note, however, that this is true only within the Israel-Stewart approach to heat conduction. In fact, the ``second-order term'' $\beta_1 \delta q^k \delta q_k$ is crucial to the correspondence, as it carries all information about the value of $\rho_S$. At very low temperatures, if the thermal excitations are phononic, the third equation of \eqref{beta1} becomes  $\beta_1 = (Tsc_{\text{ph}}^2)^{-1}$, where $c_{\text{ph}}^2$ is the speed of phonons \cite{landau6}.





% \textit{Burgers viscoelasticity is equivalent to MIS$^*$ -} Consider a fluid at zero chemical potential, whose perturbation-fields are $\varphi^A=\{\delta T,\delta u^k, \delta \Pi^{kl},\delta \Lambda^{kl}\}$, namely temperature, flow velocity, shear stresses, and an auxiliary 2tensor field. Both $\delta \Pi^{kl}$ and $\delta \Lambda^{kl}$ are symmetric and traceless. The most general information current and entropy production rate (up to field redefinitions \cite{GavassinoNonHydro2022}) are
% \begin{equation}\label{infoBurgers}
% \begin{split}
% TE^0 ={}& \dfrac{c_v (\delta T)^2}{2T} + \dfrac{1}{2}(\rho{+}P)\delta u^k \delta u_k \\
% & +\dfrac{1}{2} \beta_2 (\delta \Pi^{kl}\delta \Pi_{kl}+\delta \Lambda^{kl}\delta \Lambda_{kl})\, , \\
% TE^j ={}& s  \delta T \delta u^j + \delta \Pi^{jk}\delta u_k \, , \\
% T\sigma ={}& \xi_1 \delta \Pi^{kl}\delta \Pi_{kl}+2\xi_2 \delta \Pi^{kl}\delta \Lambda_{kl}+\xi_3\delta \Lambda^{kl}\delta \Lambda_{kl} \, . \\
% \end{split}
% \end{equation}
% If we write the field equations \eqref{cinque} explicitly, we recover the usual equations of a shear-viscous fluid in the Landau frame \cite{OlsonLifsh1990}, with the only difference that, if $\xi_2 \neq 0$, the viscous stress $\delta \Pi^{kl}$ does not obey a relaxation equation as in the Israel-Stewart theory. Instead, it obeys a second-order equation,  
% \begin{equation}\label{burgers}
% (\lambda_2 \partial_t^2  + \lambda_1 \partial_t + 1) \delta \Pi_{kl} = - 2(\eta+\chi \partial_t) \braket{\partial_k \delta u_l} ,
% \end{equation}
% where $\braket{...}$ denotes the symmetric traceless part and
% \begin{equation}
% \begin{bmatrix}
% \lambda_2 \\
% \lambda_1 \\
% 2\eta \\
% 2\chi \\
% \end{bmatrix}
% = \dfrac{1}{\xi_1 \xi_3 - \xi_2^2} 
% \begin{bmatrix}
% \beta_2^2 \\
% \beta_2 (\xi_1+\xi_3)\\
% \xi_3 \\
%  \beta_2 \\
% \end{bmatrix} .
% \end{equation}
% In non-relativistic rheology, equation \eqref{burgers} is known as ``Burgers model'' of viscoelasticity \cite{Findley_book,Prieto2014,Malek2018}. It describes a material that behaves like a fluid at low frequencies, with shear viscosity $\eta$, and like a solid at high frequencies, with shear modulus $G=\chi/\lambda_2=(2\beta_2)^{-1}$. The main difference with the Israel-Stewart theory is that the Burgers model presents two relaxation times, instead of just one. 

% If we introduce the notation
% \begin{equation}
% \begin{split}
%  z={}& \text{arcsinh}\bigg( \dfrac{\xi_3-\xi_1}{2\xi_2} \bigg),\\
% \delta \Pi^{kl} ={}& \delta \Pi^{kl}_{(1)} + \delta \Pi^{kl}_{(2)} \, , \\
% \delta \Lambda^{kl} ={}& e^z \, \delta \Pi^{kl}_{(1)} - e^{-z} \, \delta \Pi^{kl}_{(2)} \, , \\
% 2\eta_{(1)} ={}& (\xi_1+2\xi_2 e^z + \xi_3 e^{2z})^{-1}, \\
% 2\eta_{(2)} ={}& (\xi_1-2\xi_2 e^{-z} + \xi_3 e^{-2z})^{-1}, \\
% \tau_{(1)}={}& 2\eta_{(1)}\beta_2 (1+e^{2z}) \, , \\
% \tau_{(2)} ={}& 2\eta_{(2)}\beta_2 (1+e^{-2z})\, , \\
% \end{split}
% \end{equation} 
% equation \eqref{infoBurgers} becomes
% \begin{equation}\label{infoKe}
% \begin{split}
% TE^0 ={}& \dfrac{c_v (\delta T)^2}{2T} + \dfrac{1}{2}(\rho{+}P)\delta u^k \delta u_k \\
% & +\dfrac{\tau_{(1)}}{4 \eta_{(1)}} \delta \Pi^{kl}_{(1)}\delta \Pi_{(1)kl}+\dfrac{\tau_{(2)}}{4 \eta_{(2)}} \delta \Pi^{kl}_{(2)}\delta \Pi_{(2)kl}\, , \\
% TE^j ={}& s  \delta T \delta u^j + \big(\delta \Pi^{jk}_{(1)}+\delta \Pi^{jk}_{(2)} \big)\delta u_k \, , \\
% T\sigma ={}&  \delta \Pi^{kl}_{(1)}\delta \Pi_{(1)kl}/2\eta_{(1)}+\delta \Pi^{kl}_{(2)}\delta \Pi_{(2)kl}/2\eta_{(2)} \, . \\
% \end{split}
% \end{equation}
% These are the information current and entropy production rate of the so-called MIS$^*$ theory, formulated by \citet{KeYin2022} as a model of the quark-gluon plasma in the extended hydrodynamic regime. Indeed, if we write \eqref{cinque} explicitly, we recover the linearised field equations of MIS$^*$, expressed in a symmetric-hyperbolic form. In a nutshell, MIS$^*$ models a shear-viscous fluid, whose shear stress tensor is the sum of two parts, $\delta \Pi_{(1)}^{kl}$ and $\delta \Pi^{kl}_{(2)}$, subject to two independent relaxation equations:
% \begin{equation}\label{KeYinequations}
% \begin{split}
% (\tau_{(1)}\partial_t + 1)\delta \Pi_{(1)kl} ={}& -2\eta_{(1)} \braket{\partial_k \delta u_l} \, , \\
% (\tau_{(2)}\partial_t + 1)\delta \Pi_{(2)kl} ={}& -2\eta_{(2)} \braket{\partial_k \delta u_l} \, .\\
% \end{split}
% \end{equation}
% By Theorem \ref{theo2}, we can conclude that the Burgers model and MIS$^*$ are mathematically equivalent (in the linear regime). This implies, for example, that MIS$^*$ features an elastic behavior at high frequencies. More importantly, it implies that the Burgers equation \eqref{burgers} can \textit{always} be decomposed in the form \eqref{KeYinequations}, independently from whether we are able to assign a microscopic meaning to the individual contributions $\delta \Pi_{(1)}^{kl}$ and $\delta \Pi^{kl}_{(2)}$. In fact, equation \eqref{KeYinequations} can be viewed just as a convenient reparameterization of the Burgers degrees of freedom. The four parameters in \eqref{burgers} can be expressed in terms of the parameters in \eqref{KeYinequations} as follows:
% \begin{equation}\label{ultimatum}
% \begin{split}
% \lambda_2 ={}& \tau_{(1)}\tau_{(2)}  \, ,\\
% \lambda_1 ={}& \tau_{(1)}+\tau_{(2)} \, ,  \\
% \eta ={}& \eta_{(1)}+\eta_{(2)} \, ,  \\
% \chi ={}& \eta_{(1)} \tau_{(2)}+\eta_{(2)}\tau_{(1)} \, . \\
% \end{split}
% \end{equation}
% If we plug \eqref{ultimatum} into \eqref{burgers}, and we introduce the partial shear moduli $G_{(1)}=\eta_{(1)}/\tau_{(1)}$ and $G_{(2)}=\eta_{(2)}/\tau_{(2)}$ \cite{BaggioliHolography}, we obtain
% \begin{equation}
% \begin{split}
% & \dfrac{\eta_{(1)}\eta_{(2)}}{G_{(1)}G_{(2)}} \partial_t^2 \delta \Pi_{kl} + \bigg(\dfrac{\eta_{(1)}}{G_{(1)}} +\dfrac{\eta_{(2)}}{G_{(2)}} \bigg) \partial_t \delta \Pi_{kl} + \delta \Pi_{kl} \\ 
% & {=} {-} 2(\eta_{(1)}{+}\eta_{(2)}) \! \braket{\partial_k \delta u_l} {-}2 \dfrac{\eta_{(1)}\eta_{(2)} (G_{(1)}\!{+}G_{(2)})}{G_{(1)}G_{(2)}} \partial_t \! \braket{\partial_k \delta u_l} \! ,\\
% \end{split}
% \end{equation}
% which in non-relativistic rheology is a well-known  ``representation'' of the Burgers model \cite{Malek2018}.